% !TEX TS-program = pdflatexmk
\documentclass{article}

\usepackage[eandd,final]{neurips_2026}

\usepackage[utf8]{inputenc}
\usepackage[T1]{fontenc}
\usepackage{hyperref}
\usepackage{url}
\usepackage{booktabs}
\usepackage{amsfonts}
\usepackage{amssymb}
\usepackage{nicefrac}
\usepackage{microtype}
\usepackage{graphicx}
\usepackage{subcaption}
\usepackage{amsmath}
\usepackage{multirow}
\usepackage{xcolor}
\usepackage{colortbl}
\usepackage{enumitem}

\makeatletter
\newcommand{\applabel}[2]{\protected@edef\@currentlabel{#1}\label{#2}}
\makeatother

\title{Cosine is Nearly Human: The Reproducibility Ceiling of Perceptual Similarity}

\author{%
  Benjamin TenWolde \\
  Hal{\i}c{\i}o\u{g}lu Data Science Institute\\
  University of California, San Diego\\
  \texttt{betenwolde@ucsd.edu} \\
  \And
  Virginia R. de Sa \\
  HDSI and Cognitive Science\\
  University of California, San Diego\\
  \texttt{desa@ucsd.edu}
}

\begin{document}

\maketitle

\begin{abstract}
The NIGHTS dataset is widely used to evaluate whether image-similarity metrics match human perceptual judgments for mid-level equivalence of features. DreamSim, the model introduced with NIGHTS, broke ground with reports of  large gains over DINO-based baselines and has been taken as evidence that supervised training on NIGHTS produces more human-aligned visual representations. We argue that this conclusion is not supported by the benchmark as presented.  NIGHTS evaluation is constructed from the unanimous-agreement subset of a larger 100{,}000-triplet corpus, excluding a significant fraction of weak-majority and split-vote triplets where human judgments are not in agreement. We re-evaluate twelve perceptual-similarity models on this larger released corpus. On the 64{,}420 non-tie triplets, frozen DINO backbone features with plain cosine distance reach 74.1\% accuracy without any training. However, the vote counts bound what the labels can certify: a fresh panel of the same size, drawn by resampling each triplet's votes, would reproduce the observed majority only about 89\% of the time overall, and only 74--84\% on the weak-majority triplets that the public release discards. We note that the reported increase in accuracy reflects the fitting of the collected vote sample. We claim that this is not evidence that the metric would better match a new set of human judgments. We note that the source of their overly positive conclusion is in the benchmark design, not typical overfitting. NIGHTS collapses multi-axis perceptual judgments into a single binary choice, and its two human labels conflict on the same images: JND voters judged 82\% of unanimous 2AFC choices different from the reference. Under rotation, color, blur, and foreground/background ablations, we find that fine-tuning also increases model fragility, suggesting that higher in-distribution accuracy can come at the cost of less stable behavior. No human judgments exist for these manipulated stimuli, so we report this as a property of the metric, not as a model--human divergence. The pattern does not transfer to low-level distortions: on BAPPS, LPIPS beats every zero-shot DINO backbone on color, deblur, and super-resolution, so plain cosine is not competitive for low-level distortion detection. We conclude that NIGHTS supports a narrower claim than previously assumed: cosine distance on self-supervised features captures much of the reproducible perceptual signal available in the current labels. Stronger claims about human alignment require a redesigned benchmark with more votes per triplet, axis-specific similarity questions, per-subject or per-rater statistics, and robustness ablations.
\end{abstract}

\section{Introduction}

We make the claim that plain cosine distance on a self-supervised DINO-based backbone captures most of the human perceptual-similarity signal currently measurable in the image-triplet judgments from NIGHTS~\citep{fu2023dreamsim}. The competing claim is that supervised triplet training on NIGHTS (DreamSim and successors) yields strictly more human-aligned representations: DreamSim, a LoRA-fine-tuned DINO+CLIP+OpenCLIP ensemble trained on NIGHTS, reports 96.2\% on the public 20{,}019-triplet release and subsequent work targets that number. We argue that the cosine baseline gives the narrower interpretation because NIGHTS, the benchmark that should adjudicate whether supervised triplet training improves human alignment, is itself filtered. The 20{,}019-triplet release is the unanimous-agreement subset of a 100{,}000-triplet collection; 35{,}580 ties (50/50 splits after aggregation) and roughly 44{,}000 weak-majority triplets (patterns such as 6-1, 5-1, 3-1, 2-1) are discarded. We note that human annotators disagreed on many triplets, and no further analysis or evaluation was published on these disagreements. Annotator-level data is also missing, and the benchmark's binary labels do not distinguish contested triplets from unanimous ones. We argue that the cosine baseline (such as provided with base DINO) is therefore not just a strong starting point, but it supplies most of the signal these labels can express. No stronger alignment claim trained against the same labels can be certified on them. The remaining gap is a labeling problem, and the successor benchmark required to move past it is what this paper motivates.

This is, to our knowledge, the first evaluation on the full NIGHTS corpus. We evaluate the 64{,}420 non-tie triplets by 5-fold cross-validation, training on 51{,}536 and testing on 12{,}884 per fold. We exclude all 35{,}580 ties from training and hold them out to audit bias: 35{,}569 are 1-1 and 11 are 2-2, since voting stopped at the first disagreement. NIGHTS provides only vote tallies, so we treat votes as anonymous ballots and resample each triplet's own votes to ask how often a fresh panel of the same size would reproduce its majority. This gives a reproducibility ceiling that accounts for the noise in the small per-triplet vote counts (47\% of non-tie triplets have only 3 or 4 voters). We stratify test accuracy by vote pattern (X-0, X-1, X-2) and check left- versus right-candidate preference on tied votes.

Three findings support this claim.  (i) Resampling the per-triplet votes gives a closed-form reproducibility ceiling for any 2AFC benchmark: on NIGHTS a fresh panel of the same size reproduces the observed majority about 89\% of the time overall, and only 74--84\% on the 2-1 and 3-1 triplets that make up 46\% of the corpus. (ii) Cosine on the CLS embedding of a frozen DINOv1 backbone reaches 74\% with no training. LoRA on NIGHTS-20k raises clean accuracy by 4--6 percentage points (pp) but increases foreground-only fragility from 2.0 to 7.7 pp. The tuned models rely on background cues humans do not need to recognize objects~\citep{xiao2020backgrounds}. (iii) The two NIGHTS labels, 2AFC (``which is more similar'') and JND (``same or different''), conflict on the same image pairs: JND voters judged 82\% of unanimous 2AFC winners different from the reference (94\% when they too were unanimous), so a unanimous 2AFC label means ``closer than the alternative,'' not ``similar,'' and neither label is a fixed ground truth. These results show that cognitive heads (Tversky, GCM, TNN) and fine-tuning do not beat plain cosine by more than the label noise can certify. Plain cosine on DINO-family backbones is the strongest zero-shot model. The appendix lists full per-fold numbers, additional ablations (harmonization, C-RADIOv4 distillation, Qwen3-VL-as-judge), BAPPS cross-benchmark sanity checks, paired significance tests, ImageNet/non-ImageNet splits, and qualitative triplet traces (Appendix~\ref{app:traces}).

\section{Related work}

LPIPS~\citep{zhang2018lpips}, generalizing pixel-domain SSIM~\citep{wang2004ssim}, established deep-feature cosine distance as a perceptual metric. Subsequent work extended this with adversarially certified~\citep{lipsim2024} and adversarially trained~\citep{croce2025robustclip} variants. DreamSim~\citep{fu2023dreamsim} introduced synthetic Stable Diffusion v1.4 (SD)~\citep{rombach2022stablediffusion} triplets and LoRA fine-tuning~\citep{hu2021lora} for human alignment, reporting 96.2\% on the filtered 20k and becoming the default backbone for follow-up work including DiffSim~\citep{song2025diffsim} and the perceptual-alignment transfer analysis of \citet{sundaram2024perceptual}. All of these evaluations use the same 20k filtered split as the original DreamSim paper.

Self-supervised backbones are the strongest perceptual base. \citet{caron2021dino} showed that DINO's self-attention recovers object structure without supervision. \citet{muttenthaler2022vice}, \citet{muttenthaler2022human}, and \citet{conwell2024large} found that SSL ViTs align with human similarity and brain responses above supervised CNNs, complementing earlier evidence that ImageNet~\citep{deng2009imagenet} accuracy alone is a poor predictor of perceptual similarity~\citep{kumar2022perceptual}. \citet{oquab2023dinov2} scaled this to DINOv2 and DINOv3, and \citet{chen2025midlevel} benchmarked 22 SSL backbones on mid-level vision tasks including NIGHTS. A second line of work injects human alignment at training time: AligNet~\citep{muttenthaler2025alignet} learns an affine transform into a THINGS-derived human similarity space and distills it into vision backbones, and UniSim~\citep{ghazanfari2025unisim} trains a multi-task metric over 25 datasets including NIGHTS. A third line argues that attribute-similarity metrics (LPIPS, CLIP~\citep{radford2021clip}, DINO, DreamSim) miss \emph{relational} similarity, the shared logic between visually different scenes, and proposes a VLM-based metric for this axis~\citep{le2025relsim}. Across all of these approaches the final similarity operator is cosine distance (or a dot product on unit-norm embeddings), as is the case in retrieval and recommendation systems where cosine dominates despite its simplicity~\citep{steck2024cosine,you2025semantics}. We do not propose a new alignment method. We audit what NIGHTS can and cannot tell us about any model's human alignment, and we show that cosine distance captures most of the reliable agreement in NIGHTS labels.

Cognitive theories of similarity provide alternative inductive biases. Tversky's~\citep{tversky1977} contrast model uses asymmetric feature overlap in place of distance. Nosofsky's generalized context model adds selective dimensional attention~\citep{nosofsky1986}. Gentner's structure-mapping theory aligns relational representations~\citep{gentner1983}. \citet{doumbouya2025tversky} recently introduced a fully differentiable Tversky Neural Network. We evaluated these cognitive architectures directly against plain cosine-distance heads.

\section{Method}

\subsection{Dataset reconstruction}

The public NIGHTS release contains 20{,}019 triplets with binary vote columns. This is the unanimous subset of a larger 100{,}000-triplet collection that the authors released with raw vote proportions and per-triplet vote counts. The authors confirm that participant-level data was not retained, so we could not analyze individual variability across the 100{,}000 triplets. We converted the proportions to integer counts, computed the agreement pattern (larger-smaller, e.g., ``6-1'' for 6 left and 1 right), and partitioned the corpus into 64{,}420 non-tie triplets for training and evaluation and 35{,}580 tie triplets held out for the bias audit. DreamSim's trained and published results are on the  20k subset. This is the union of the X-0 (unanimous) patterns in the non-tie set. Its 20{,}019 triplets include one with a single vote and two with eleven; analyses that restrict the unanimous subset to 3--10 votes, or drop one of those extremes, report $N=20{,}005$, $20{,}017$, or $20{,}018$. The vote-count distribution is skewed to small counts: 28.9\% of non-tie triplets have 3 votes, 18.2\% have 4, and 5.5\% have at least 9.

\subsection{Evaluation protocol}

We use five-fold cross-validation with the commonly used seed 42 on the 64,420 non-tie triplets. Each fold has 12,884 test triplets and 51,536 training triplets. For LoRA, we hold out 10\% of training for validation. We report mean and standard deviation of test accuracy across five folds by agreement pattern, and predictions on each fold's disjoint 7,116-triplet share of the 35,580 held-out ties.


\paragraph{Reproducibility ceiling.} We estimate how often an independent panel of voters would reproduce a triplet's majority label by resampling its observed votes, assuming exchangeable voters (\S\ref{sec:limitations}, vi). For each non-tie triplet with $L$ left votes and $R$ right votes ($V = L+R$), we draw $V$ votes with replacement from the observed proportions and ask whether the new panel's majority matches the observed one; a tie in the new panel counts half. The binomial form gives this probability in closed form per $(L,R)$ pattern. Averaging over 64{,}420 triplets gives the overall ceiling. This is an upper bound on accuracy these labels can certify: agreement beyond it may fit votes that would not recur. It does not bound accuracy against the fixed observed labels, which is trivially 100\%.

\subsection{Models evaluated}

A model is a backbone, an adaptation (frozen, LoRA, retrained, or released weights), and a head that turns embeddings into a similarity; names follow the template \texttt{\{Backbone\}-\{Adaptation\}-\{Head\}-\{Dataset\}}. Table~\ref{tab:naming} (appendix) lists the factor values. Accuracies are on NIGHTS-100k unless stated otherwise. For zero-shot models we used cosine distance on the post-norm CLS token at 224$\times$224 input. For LoRA models we applied rank-8 LoRA to the query, key, value, and output projections of every transformer block, unfroze the final LayerNorm, and trained with Schedule-Free AdamW~\citep{defazio2024schedulefree} (hyperparameters in Appendix~\ref{app:training}). The four head variants share the same LoRA budget: plain cosine distance, our 197k-parameter Cog head (frozen cosine alignment, GCM context, and Tversky contrast with learnable $\alpha, \beta$), the inverse-frequency-stratified Cog, and the 197k-parameter Tversky Neural Network of \citet{doumbouya2025tversky} (product intersection and subtract-match difference functions, 256 features). DreamSim-Pretrained used the authors' public weights. DreamSim-Scratch retrained the full 3-model ensemble with rank-16 LoRA and margin-0.05 hinge loss for 6 epochs on each fold's training set using the published codebase unchanged, giving one checkpoint per fold.

\section{Results}

\subsection{Main accuracies and the cost of unfiltering}

Table~\ref{tab:main} reports 5-fold mean accuracy on the full NIGHTS-100k for all primary models. Per-fold numbers are in Table~\ref{tab:perfold} of the appendix. Figure~\ref{fig:drop} shows each model's accuracy on the unanimous-only X-0 subset (comparable to the published 20k) against its accuracy on the 100k corpus. We found that every model lost between 10.6 pp (LPIPS) and 17.2 pp (DreamSim-Pretrained) when evaluated on the full corpus vs the unanimous 20k published subset. On unanimous triplets the strongest models were near ceiling, while on weak-majority triplets they were near chance. NIGHTS-20k kept only the unanimous triplets, so the published scores were measured on the easy part of the data. Figure~\ref{fig:drop} shows the result as increasing disagreement is reincorporated. 

\begin{table}[ht]
\caption{5-fold CV accuracy for selecting the human-majority candidate on NIGHTS-100k non-tie ($N=12{,}884$ per fold). $^\dagger$DreamSim-Pretrained uses released weights trained on the NIGHTS-20k training split; some of those triplets occur in our test folds, inflating its score. DreamSim-Scratch is retrained from scratch on NIGHTS-100k using the authors' published training recipe unchanged.}
\label{tab:main}
\centering
\small
\begin{tabular}{llcc}
\toprule
Model & Backbone + Head & Accuracy & Trainable params \\
\midrule
LPIPS-AlexNet-Pretrained & AlexNet + linear probe & $65.3 \pm 0.2$ & 0 \\
DINOv1-ZS-CLS-cos & DINOv1 ViT-B/16 & $\mathbf{74.1 \pm 0.5}$ & 0 \\
DINOv2-ZS-CLS-cos & DINOv2 ViT-B/14 & $70.6 \pm 0.4$ & 0 \\
DINOv3-ZS-CLS-cos & DINOv3 ViT-B/16 & $72.0 \pm 0.4$ & 0 \\
\midrule
DINOv1-LoRA-Cog & DINOv1 + Cog & $75.1 \pm 0.3$ & 641{,}346 \\
DINOv2-LoRA-Cog & DINOv2 + Cog & $77.1 \pm 0.5$ & 641{,}346 \\
DINOv2-LoRA-Cos & DINOv2 + Cosine & $78.4 \pm 0.4$ & 443{,}904 \\
DINOv2-LoRA-Cog-strat & DINOv2 + Cog + strat & $76.3 \pm 0.6$ & 641{,}346 \\
DINOv2-LoRA-TNN & DINOv2 + Doumbouya 2025 TNN & $77.4 \pm 0.5$ & 640{,}515 \\
DINOv3-LoRA-Cog & DINOv3 + Cog & $\mathbf{78.7 \pm 0.3}$ & 641{,}346 \\
\midrule
DreamSim-Pretrained$^\dagger$ & DINO+CLIP+OpenCLIP, 20k weights & $80.7 \pm 0.2$ & $\sim$1.8M \\
DreamSim-Scratch & Retrained on NIGHTS-100k, votes$\geq$6 & $78.7 \pm 0.7$ & $\sim$1.8M \\
\bottomrule
\end{tabular}
\end{table}

\begin{figure}[ht]
\centering
\includegraphics[width=0.78\linewidth]{figures/drop_chart.pdf}
\caption{The same models use 5-fold cross-validation on unanimous X-0 test folds (blue; disjoint from training, comparable to published 20k) and all non-tie NIGHTS-100k triplets (red). Each drops 10 to 17 points because test sets differ. Dashed line: DreamSim's published 96.2\% on their 20k held-out split.}
\label{fig:drop}
\end{figure}

\subsection{Agreement-stratified reliability}

Figure~\ref{fig:agreement} plots accuracy against the vote-majority fraction for each model. The curves rise with agreement and share one shape: near ceiling on unanimous triplets, falling as disagreement grows, and close to chance on 2-1 triplets, where the human majority is a single vote. The spread between models widens with agreement: on 2-1 triplets the worst (LPIPS, 57.0\%) and the best (DreamSim, 65.2\%) differ by 8.2 pp, on 3-1 triplets the full spread is 13 pp, and on unanimous triplets it is about 20 pp. In other words, the more accurate models pull ahead only where human agreement is already strong; where people were split, all models sit near chance together. The extra accuracy comes from the triplets whose labels are reliable, not from doing better on the ones where people disagree.
\begin{figure}[ht]
\centering
\includegraphics[width=0.78\linewidth]{figures/agreement_curve.pdf}
\caption{Accuracy versus human-agreement strength on NIGHTS-100k non-tie triplets. Accuracy rises with agreement for every model; DreamSim leads at most levels, but its released weights overlap our test folds. Unanimous patterns are pooled at 1.0; patterns with fewer than 100 triplets are omitted. Shaded band: the unanimous NIGHTS-20k subset.}
\label{fig:agreement}
\end{figure}

\subsection{A reproducibility ceiling from the vote counts}
\label{sec:ceiling}

Table~\ref{tab:ceiling} reports the reproducibility ceiling for each subset: the probability that a fresh panel of the same size, drawn by resampling a triplet's observed votes with replacement, reproduces the observed majority label (a tie in the new panel counts half). The assumption is that each triplet's observed votes are exchangeable draws from one voter population (\S\ref{sec:limitations}, vi). The ceiling is a property of the labels, not of any model. We report it as a ceiling on the accuracy these labels can certify, not a bound on accuracy against the fixed observed labels: a deterministic model that memorized the observed majorities reaches 100\%. The question is what a model generalizing from those labels can be expected to reproduce on a fresh vote. The overall ceiling on NIGHTS-100k non-tie is about 89\%. The aggregate hides a split: 100\% on unanimous triplets but far less where the majority hangs on one or two votes: a fresh panel of three agrees with a 2-1 majority only 74\% of the time (20/27; $N{=}18{,}454$), a fresh panel of four with a 3-1 majority 84\% ($N{=}11{,}181$), and a fresh panel of five with a 4-1 majority 94\% ($N{=}6{,}841$). The 2-1 and 3-1 patterns together account for 46\% of the non-tie corpus and drag the ceiling down.

\begin{table}[ht]
\caption{Reproducibility ceiling on NIGHTS-100k: the probability that a fresh panel of the same size, drawn by resampling each triplet's observed votes with replacement, reproduces the observed majority label. The strong-majority row spans its patterns (5-1 through 10-1); the overall value lies between 88.8\% and 89.2\% for any mix.}
\label{tab:ceiling}
\centering
\begin{tabular}{lcc}
\toprule
Subset & N & Repeatability \\
\midrule
Unanimous (X-0) & 20{,}018 & 100.0\% \\
Strong majority (X-1, $X \geq 5$) & 7{,}922 & 96.5--99.5\% \\
Weak majority (4-1, 3-1, 2-1) & 36{,}476 & 81.0\% \\
\midrule
All non-tie (NIGHTS-100k) & 64{,}420 & $\approx$89\% \\
\bottomrule
\end{tabular}
\end{table}

No model reaches this ceiling: the best, DreamSim-Pretrained, scores 80.7\% against the $\approx$89\% the labels could certify, and on 2-1 triplets, where the label itself repeats only 74\% of the time, every model sits below that bar (57.0--65.2\%, Fig.~\ref{fig:agreement}). What the ceiling rules out is reading small gaps as better alignment: differences of a few percentage points fall inside the label noise, so they cannot be attributed to one model matching people better. The models that gained the most clean-image accuracy also lost invariance (\S\ref{sec:shape_axes}), and 96.6\% of DreamSim-Pretrained's disagreements with the majority label matched at least one dissenting voter.

The root cause is the NIGHTS collection procedure, which solicited votes until the first disagreement and then stopped. As a result, 47\% of non-tie triplets were labeled by 3 or 4 voters. For a triplet with true preference $p = 0.6$, the probability that a 3-voter sample produces a majority in the other direction is $\binom{3}{0}(0.6)^0(0.4)^3 + \binom{3}{1}(0.6)^1(0.4)^2 = 0.352$, so better than one in three observed majority labels would flip under re-annotation. A stable majority at $p=0.6$ requires on the order of 35 votes per triplet. Model differences of a few percentage points on NIGHTS-100k fall within this label-noise band. They cannot be attributed to model quality without re-annotation data.

\subsection{Cognitive priors and backbone ordering}

We asked whether a psychologically motivated similarity head beats plain cosine. We tested three, Tversky asymmetric contrast, GCM selective attention, and the differentiable TNN of \citet{doumbouya2025tversky}. None did: cosine reached 78.4\% $\pm$ 0.4 against 77.1\% for Cog, 76.3\% for inverse-frequency stratification, and 77.4\% for TNN. The Tversky head shows why. It has two free weights, $\alpha$ for features only the reference has and $\beta$ for features only the candidate has; unequal values would make similarity directional, but the fitted values came out nearly equal, so the head reduced to a symmetric measure and behaved like cosine. Retraining DreamSim from scratch on NIGHTS-100k with the authors' code gave $78.7 \pm 0.7$\%, 2.0 pp below its published weights on the same splits. Both gaps are smaller than the label noise: on 2-1 and 3-1 triplets a fresh panel would agree with the observed majority only 74--84\% of the time, and differences under $\sim$10 pp are not separable from that noise (Section~\ref{sec:ceiling}). A Cohen's-$\kappa$ inter-model analysis is in Appendix~\ref{app:kappa}.

On NIGHTS-100k, zero-shot DINOv1 (74.1\%) beat DINOv2 (70.6\%) and DINOv3 (72.0\%). Under LoRA the order reversed: DINOv3-LoRA-Cog 78.7\%, DINOv2-LoRA-Cos 78.4\%, DINOv1-LoRA-Cog 75.1\%. Zero-shot rank does not predict fine-tuned rank (DINOv3 gained 6.7 pp).

\paragraph{Backbone family and scale.} Cosine is not DINO-specific. We ran sixteen zero-shot backbones: eight CNNs (five supervised, three bio-inspired), two non-DINO ViTs (CLIP and supervised), and six DINO ViTs (DINOv1, DINOv3, and DINOv2 at S, B, L, and g). The ten non-DINO models span 65.3--72.4\%, and six of them score above 70\% (Table~\ref{tab:backbones}). Bigger is not better: DINOv2 falls monotonically over 52$\times$ in size on matched folds, and its two largest fall more than 5 pp behind DINOv1 cosine (Table~\ref{tab:scale}), extending the CLIP/MAE trend of \citet{fu2023dreamsim}. Recent large models do worse: C-RADIOv4-SO400M at 58.2\% after one fold-0 epoch and Qwen3-VL-4B at 63.0\% (Appendix~\ref{app:extra}).

\begin{table}[ht]
\caption{Zero-shot DINOv2 with cosine on NIGHTS-100k (5-fold mean $\pm$ s.d.). Accuracy falls monotonically with size. The 70.1\% reference is the split-half agreement of the NIGHTS votes, the probability that the two halves of a triplet's votes choose the same majority; it is a conservative companion to the $\approx$89\% fresh-panel ceiling of \S\ref{sec:ceiling}.}
\label{tab:scale}
\centering
\small
\begin{tabular}{lcccc}
\toprule
DINOv2 & Params & NIGHTS-100k & vs.\ 70.1\% & Unanimous X-0 \\
\midrule
ViT-S/14 & 21M  & $72.2 \pm 0.5$ & $+2.1$ & 85.8 \\
ViT-B/14 & 86M  & $70.6 \pm 0.4$ & $+0.5$ & 83.9 \\
ViT-L/14 & 300M & $69.0 \pm 0.2$ & $-1.1$ & 81.8 \\
ViT-g/14 & 1.1B & $68.4 \pm 0.3$ & $-1.7$ & 80.5 \\
\bottomrule
\end{tabular}
\end{table}

\paragraph{The fine-tuning gain is real but not certifiable as alignment.}\label{sec:gain} The untrained DreamSim ensemble reached 74.5\% versus 74.1\% for frozen DINOv1 (0.4 pp from ensembling). Retraining reached 78.7\% (4.2 pp from training), while the pretrained checkpoint gained 2.0 pp overall and 3.8 pp on its overlapping 20k training subset. The 4.2 pp gain is not a color artifact: splitting by candidate color differences yields similar gains ($+3.9$ and $+4.9$ pp; Appendix~\ref{app:colorcontrol}). Training learns real signal from the labels, but their 74--84\% weak-majority reproducibility cannot establish better human agreement (\S\ref{sec:ceiling}).

\subsection{The ceiling and the ranking on a second benchmark (BAPPS)}
\label{sec:bapps}

The same resampling applies to BAPPS~\citep{zhang2018lpips}, which collects 5 votes per triplet ($N=36{,}344$). A 3-2 majority is reproduced by a fresh panel of five only 68.3\% of the time and a 4-1 majority 94.2\%, so 2AFC benchmarks with few votes per triplet carry the same per-pattern limits on what their labels can certify.

On BAPPS, the ranking reverses: LPIPS beats every zero-shot DINO backbone by 1.4 to 5.6 pp, while DINO leads on NIGHTS by 5 to 10 pp. Only DreamSim-Pretrained (BAPPS 73.6, NIGHTS 80.7) and zero-shot DINOv1 (71.5, 74.1) are competitive on both (Table~\ref{tab:bapps}). Plain cosine is not competitive for low-level distortions; no single score captures performance across both benchmarks. Per-category comparisons appear in Appendix~\ref{app:crossover}.

\begin{table}[ht]
\caption{Zero-shot 2AFC accuracy on BAPPS validation (six distortion categories) and NIGHTS-100k. LPIPS leads on BAPPS, DINOv2/v3 on NIGHTS; only DreamSim-Pretrained and DINOv1 are competitive on both.}
\label{tab:bapps}
\centering
\small
\begin{tabular}{lccc}
\toprule
Model & BAPPS & NIGHTS-100k & Above 69.0\,/\,70.1\% \\
\midrule
LPIPS-AlexNet           & 72.9 & 65.3 & \\
DINOv1-ZS               & 71.5 & 74.1 & \checkmark \\
DINOv2-ZS               & 67.3 & 70.6 & \\
DINOv3-ZS               & 68.4 & 72.0 & \\
DreamSim-Pretrained & 73.6 & 80.7 & \checkmark \\
\bottomrule
\end{tabular}
\end{table}

\subsection{Shape, texture, and the axes of perceptual alignment}
\label{sec:shape_axes}

Accuracy against NIGHTS majority labels shows how often a model agrees, not which image features it uses. We therefore replicate the DreamSim appendix's stimulus-manipulation study~\citep{fu2023dreamsim}, measuring accuracy under eleven manipulations.

\paragraph{Eleven stimulus manipulations.} We apply each manipulation identically to all three images in a triplet from the unanimous 20k subset ($N=2000$ sampled for speed), then score accuracy against the unchanged human label. The eleven manipulations cover four of the five DreamSim Figure 6 axes: orientation (90/180/270$^\circ$), color (grayscale), shading (HSV V-channel histogram equalization), structure (Gaussian blur $\sigma\in\{1,2,4,8\}$), and foreground/background (center disc of radius 0.35$\cdot$image side or its complement). NIGHTS has no human labels for manipulated images, and we collected none. These tests measure model invariance, not human perceptual agreement; per-axis results are in Appendix~\ref{app:axes_detail}.

\paragraph{The same backbone, trained longer, loses invariance.} Across 16 backbones, higher clean accuracy correlates with larger changes under manipulation, but this does not isolate training. We train DINOv2 with rank-8 LoRA on NIGHTS-20k for 8 epochs and test checkpoints \{0,1,2,3,5,8\} on the same 2000 triplets. Fragility is the mean absolute change in accuracy across the eleven manipulations, using absolute values because some increase accuracy. From epoch 0 to 1, clean accuracy rises from 84.6 to 89.8\% and fragility from 1.9 to 5.1 pp. By epoch 8, accuracy falls to 85.2\% while fragility climbs to 12.3 pp. With backbone and data fixed, training duration causes the invariance loss. The trajectory and attention overlays are in Appendix~\ref{app:extra}.

\paragraph{Foreground-only fragility rises after LoRA.} NIGHTS uses Stable Diffusion images with centered subjects and low-detail, nonuniform backgrounds. The foreground-only mask keeps a centered disc of radius 0.35 times the image side (about 38\% of its area) and sets the surround to black. We apply the same mask to all three images without cropping or rescaling, preserving subject size and position. Removing the surround tests whether the model relies on its cues despite their low detail. Zero-shot DINOv2 drops 2.0 pp; DINOv2 + LoRA drops 7.7 pp, a 3.8$\times$ increase. No human labels exist for these manipulated triplets, so this measures model invariance, not model--human divergence. In related natural-image categorization tasks, humans lose $\sim$3 pp under IN-9 background randomization~\citep{xiao2020backgrounds} and $\sim$1 pp on ObjectNet~\citep{barbu2019objectnet}. The zero-shot drop is within that range; the LoRA drop is outside it, though these are not synthetic-triplet similarity judgments.

\paragraph{Fine-tuning makes models brittle; zero-shot DINO stays robust.} LoRA-trained and DreamSim-Pretrained models are both more accurate on clean images and more fragile under the 11 manipulations than every zero-shot backbone. Each point of clean accuracy gained above the zero-shot band costs about a point of mean fragility (Appendix Figure~\ref{fig:cost_acc}), while zero-shot DINOv3 is the most invariant backbone tested. What DreamSim reports as better human alignment thus comes with more brittleness than the zero-shot ViT and DINO backbones. 

\paragraph{A 2AFC lead does not carry over to JND.} No model is trained for JND; all training uses the 2AFC labels alone. NIGHTS also ships a second human task, JND, on 608 of the 2AFC triplets: each candidate was shown briefly with the reference and judged ``same'' or ``different,'' giving 1{,}216 reference--candidate pairs labeled by majority vote. We score the same frozen backbones on it, with no retraining: we reuse the embeddings from the main results and use the JND labels only to grade them, by how well the distance a backbone puts between a pair's two images separates the ``different'' pairs from the ``same'' ones (ROC AUC, which needs no cutoff). On the 2AFC triplets DINOv1 beats DINOv2 by 3.8 pp; on JND the lead disappears (AUC 0.674 vs.\ 0.677; Table~\ref{tab:axes}). Across the fourteen zero-shot models, 2AFC accuracy barely predicts JND AUC (Spearman $\rho = 0.37$, $p = 0.19$). The two models trained on NIGHTS 2AFC score highest on both tasks (DreamSim 0.727, DINOv2+LoRA 0.720), though a single AUC carries a 95\% interval of about $\pm 0.04$. One benchmark number is thus a weak stand-in for perceptual alignment (Appendix~\ref{app:axes}).

\subsection{Model confidence versus the human vote split}
\label{sec:confidence}

Accuracy only shows whether a model picks the same image as most people. It does not ask whether the model hedges when people are also split on the choice.  To check that, we turn each model's score into a probability of choosing the left image and read how confident it is (the calibration and tables are in Appendix~\ref{app:confidence}). No model is trained to output confidence; it is read off afterward with the same calibration for every model, so the differences come from the models, not the calibration. The human vote sets the target: a tie means no preference, a unanimous triplet full agreement.

The baseline is untrained cosine, and fine-tuning makes confidence worse, not better. On the 35{,}580 tied triplets, where the vote supports zero confidence, untrained DINOv1 cosine already reports 0.38 (LPIPS 0.25); fine-tuning raises this to 0.48 for DreamSim-Pretrained and to 0.80--0.95 for our own LoRA heads. Even on 2-1 triplets DreamSim-Pretrained sits at 0.49, and on weak-majority triplets, where the majority hangs on one or two votes, 11.8\% of its predictions still exceed 0.9. On unanimous triplets the same upward shift helps, since the target is high (DreamSim 0.79 against untrained DINOv1's 0.58). So the overconfidence grows with training, measured against the untrained cosine it starts from. 

Fine-tuning does have one edge, and we do not dispute it. A trained model's confidence rises and falls with the human vote more tightly than an untrained one: the two move together with rank correlation 0.76 for DreamSim-Pretrained, against 0.62 for zero-shot DINOv1 and 0.41 for LPIPS. So training teaches a model to follow how much people agree, without teaching it to be the right amount of unsure when they are split. Tables for each agreement level and the reliability curves are in Appendix~\ref{app:confidence}.

\section{Discussion}
\label{sec:discussion}

\paragraph{Why cosine is nearly human, and why more training is not quite there.} Plain cosine distance on the CLS embedding of a frozen DINO backbone reached 74\% on NIGHTS-100k with no training, and its accuracy drop under the foreground-only mask stayed inside the range external human anchors suggest. LoRA fine-tuning raised clean accuracy from 84.6 to 89.3\%, but that 4.7-pp gain came with a loss of invariance: the foreground-only drop grew from 2.0 to 7.7 pp while the object was fully visible, and DreamSim-Pretrained was the most accurate and the most fragile model across orientation, color, heavy blur, and foreground and background masks (Fig.~\ref{fig:fig6_replica}). With the same LoRA budget, no cognitive head (Tversky, GCM, TNN) beat plain cosine, and the Tversky heads learned no direction, so they reduced to symmetric measures. The 4.2 pp that training adds survives every control we ran, so it is real signal. But NIGHTS cannot tell us whether it means the model agrees better with people, or only that it has matched how one set of voters happened to weight the features in these images. The labels are too noisy to decide: a fresh panel reproduces the majority only 81\% of the time on weak-majority triplets and 74\% on 2-1, so a 5-pp gap between models falls inside that noise. The JND conflict makes the same point without any resampling (Fig.~\ref{fig:jnd_2afc_tier}). The 96.2\% DreamSim reports is measured on the 20\% of triplets where voters already agreed unanimously, where the ceiling is 100\% by construction. Retraining the same ensemble on the full 100k drops it to 94.1\% on that unanimous subset, 3.8 pp below the 97.9\% it reaches when trained on the subset alone: the published weights are tuned to exactly the easy slice the 20k release keeps. We therefore read DreamSim not as a model that matches people better, but as one fit to the specific voters and images in the 20k set. Telling genuine alignment apart from that kind of fit will take a benchmark whose labels record which feature each voter judged by, and models that compare images feature by feature instead of collapsing everything into one distance.

\paragraph{NIGHTS is a single-axis benchmark presented as a general one.} Perceptual similarity spans three levels: low-level pixel fidelity (where LPIPS and BAPPS operate), mid-level attributes of an object (pose, color, shape, texture, lighting), and high-level object identity and scene layout. NIGHTS sits between the mid and high levels, where a voter folds several mid-level attributes and a high-level judgment into one binary choice. We hypothesize that unanimous triplets differ mainly on one mid-level feature, while trade-offs across features drive the weak-majority votes in the 80{,}000 triplets NIGHTS-20k excluded. The 16-pp gap between DreamSim's 96.2\% on 20k and 80.7\% on 100k is the cost of leaving that single-axis subset. What the reproducibility ceiling exposes is not random noise. Voters disagree because they weigh different features of the same images, and the 20k release hides this by keeping only the triplets where they were in agreement. A model tuned to one benchmark inherits that benchmark's particular set of voters, so scoring well on it is not the same as matching people in general. 

\paragraph{Voting is preference, not an explanation.} Both NIGHTS 2AFC (``which is more similar?'') and NIGHTS-JND (``same or different?'') collect a single bit per voter. Multi-axis differences (shape, color, pose, lighting, composition, background, identity) collapse into one vote that records the choice but not the attribute the voter used. The two labels can be compared directly. NIGHTS-JND reuses 608 unanimous 2AFC triplets, so each of their 1{,}216 reference--candidate pairs carries a 2AFC label (its candidate won or lost) and a JND label (``same'' or ``different'' from the reference). The labels conflict when the 2AFC winner is judged ``different'' or the loser ``same.'' JND voters judged 82\% of the unanimous 2AFC winners different from the reference, 94\% when they were unanimous too, but only 6\% of the losers the same; 44\% of all pairs conflict, at every 2AFC vote count (Figure~\ref{fig:jnd_2afc_tier}). Both tasks also split internally: 47\% of JND pairs and 69\% of non-tie 2AFC triplets have a dissenting voter (Appendix~\ref{app:disagree}). A 2AFC vote is a forced comparison with no ``neither'' option, so a candidate can win unanimously by being the less different of two clearly different images. A unanimous 2AFC label thus means ``closer than the alternative,'' not ``similar to the reference.'' Whether a shared background should count toward similarity is itself one of these collapsed axes: the fragility results in \S\ref{sec:shape_axes} show that tuned models weigh it more heavily than untuned ones, but without per-rater attribute labels we cannot tell whether human voters weighed it at all. Because both labels sit on the same image pairs, the dissociation is not an image-distribution artifact. Because the tasks disagree on the same images, treating either label as ground truth chooses one meaning of ``similar''. Without per-voter labels, we cannot tell whether people perceived the images differently or interpreted the questions differently. The same release omits annotator identifiers and demographics.
\begin{figure}[ht]
\centering
\includegraphics[width=\linewidth]{figures/jnd_2afc_conflict_main.pdf}
\caption{The two human labels on the same image pairs. Left: share of unanimous 2AFC winners that JND voters judged different from the reference, and of losers judged the same, by the triplet's 2AFC vote count; 44\% of all 1{,}216 pairs conflict at every count. Right: a triplet chosen 10/10 in 2AFC whose winner (green) all three JND voters judged different from the reference. Full breakdown in Appendix~\ref{app:disagree}.}
\label{fig:jnd_2afc_tier}
\end{figure}

\subsection{A protocol for the next benchmark}
\label{sec:protocol}

Models are not the bottleneck; the labels are. Today's vision backbones already do the things perceptual similarity needs: they match points on an object as its pose changes~\citep{amir2021deepvit,caron2021dino}, stay stable when the view rotates~\citep{naseer2021intriguing,raghu2021vits}, can be tuned to judge by shape rather than texture~\citep{geirhos2018imagenettrained,hermann2020origins}, and their sensitivity to color, rotation, and blur can each be measured (Appendix~\ref{app:shape_axes_extra}). What they cannot get past is a label that records only one forced choice per triplet and discards which feature the person used to make it. A benchmark that fixed this would let us ask whether a model is right for the right reason. It would: \textbf{(1)} change one feature at a time, by taking a single image and editing it by a small and a large amount along one dimension such as brightness or blur, so the correct answer is fixed without asking anyone; \textbf{(2)} ask the same question NIGHTS asked (``which image is more like the center one?'') without naming the dimension, and after the choice have the person mark on the image what looked different, say where, and rate how sure they were, so each vote comes with its reason; \textbf{(3)} release every person's own responses, not just the vote totals, so a model's mistakes can be checked against how people actually split; \textbf{(4)} collect enough votes per triplet to push the reproducibility ceiling of \S\ref{sec:ceiling} close to 100\%, so a reported gap between models means something; \textbf{(5)} fix the train/test splits and tie rules and ship them with the data, so every model is scored the same way; and \textbf{(6)} report accuracy on held-out image distortions next to the headline number, so robustness is visible.

A first study re-annotates 3{,}000 NIGHTS triplets stratified by agreement pattern, with 30 annotators per triplet, no early stopping, and per-rater votes in the release. A second condition has the same annotators judge the \S\ref{sec:shape_axes} manipulations of those triplets, supplying the human baseline the invariance results currently lack, and the 608 triplets shared with NIGHTS-JND, so the 2AFC--JND dissociation can be traced to raters. Simply re-annotating the same images would not help: more votes would only shrink the error bars, not answer the finer questions a multi-axis test is built to ask. 

\section{Limitations}
\label{sec:limitations}

\textbf{Scope caveats.} \textbf{(i)} The model accuracies in Table~\ref{tab:main} are specific to NIGHTS, although the reproducibility ceiling itself transfers to other 2AFC benchmarks (\S\ref{sec:bapps}). \textbf{(ii)} We ran sixteen backbones zero-shot but fine-tuned only the three DINO models, and the size sweep covers one family (DINOv2, S to g); other families may scale differently. \textbf{(iii)} All LoRA models used rank 8 so that the heads were compared at equal adapter capacity; we did not sweep the rank. \textbf{(iv)} When DreamSim-Pretrained disagreed with the majority label, a minority voter had made the same choice in 96.6\% of cases; without per-voter labels we cannot tell whether those triplets were ambiguous or whether the model shares a bias with those voters. \textbf{(v)} The claim that some JND/2AFC disagreements reflect different readings of the instructions rests on one co-author looking at ten JND pairs, not a systematic re-annotation, so we treat it only as a hypothesis. \textbf{(vi)} The reproducibility ceiling resamples each triplet's observed votes assuming exchangeability, taking the observed vote fraction at face value, which is optimistic for small panels; different pools, instructions, and within-session effects would lower real agreement. Independent re-annotation reproducing the labels well above it would overturn our reading. \textbf{(vii)} Public NIGHTS votes lack rater IDs, demographics, per-rater consistency, and quality filters, so rater bias, subgroup accuracy, and majorities within a representative subpopulation cannot be tested~\citep{henrich2010weirdest}.

\section*{Broader impact and reproducibility}

We also release the seed-42 5-fold definitions, as well as  per-triplet test predictions for all models.  We release the retrained checkpoints, the standalone reproducibility-ceiling script parameterized by vote count, and the analysis notebooks that regenerate every table and figure from the released predictions, at \url{https://github.com/benlten/cosine_is_nearly_human} under the MIT license. All human judgment data were obtained from the public NIGHTS, NIGHTS-JND, and BAPPS releases.  There was no collection of new human data.

\section*{Acknowledgements and disclosures}

\textbf{Funding.} Benjamin TenWolde is supported by the Department of Defense (DoD) SMART Scholarship-for-Service Program. This work used hardware originally provided by funding from Sony and NSF IIS-CRCNS-2208362.

\textbf{Competing interests.} The authors declare no competing interests.

\bibliographystyle{plainnat}
\bibliography{references}

\appendix
\clearpage

\noindent{\Large\bfseries A\quad Method naming and factor values}
\label{app:naming}
\vspace{1em}

Throughout the paper we identify each evaluated method by a four-part name of the form \texttt{\{Backbone\}-\{Adaptation\}-\{Head\}-\{Dataset\}}. Table~\ref{tab:naming} enumerates the possible values each component can take. Backbones are the feature extractors. Adaptation describes whether the backbone is frozen, partially trained via LoRA adapters, fully retrained, or used with externally-supplied pretrained weights. The head is the module that converts two (or three) backbone outputs into a similarity score; nine variants are compared including plain cosine distance, three cognitive-science-motivated architectures (Cog-v1 through Cog), the published Tversky Neural Network of \citet{doumbouya2025tversky}, the DreamSim margin-hinge ensemble loss, and a VLM logit-ratio judge. Dataset values record which NIGHTS subset a method was trained and evaluated on; omitting the dataset suffix defaults to NIGHTS-100k, our primary evaluation target.

\begin{table}[!ht]
\caption{Method naming scheme. Every method name in the paper is a concatenation of one value from each of the four rows. Example: \texttt{DINOv3-LoRA-Cog} uses DINOv3 as backbone, rank-8 LoRA on qkv/proj, our Cog head (frozen cosine alignment + GCM context + Tversky contrast), evaluated on the full NIGHTS-100k corpus.}
\label{tab:naming}
\centering
\small
\begin{tabular}{lll}
\toprule
Component & Values & Description \\
\midrule
Backbone & DINOv1, DINOv2, DINOv3 & Self-supervised ViTs \\
 & C-RADIOv4 & Distillation of DINOv3+SigLIP2+SAM3 \\
 & Qwen3-VL-4B & Vision-language model \\
 & DreamSim-Ens & DINO+CLIP+OpenCLIP ensemble \\
 & AlexNet-lin & LPIPS classical backbone \\
\midrule
Adaptation & ZS & Zero-shot, frozen backbone \\
 & LoRA & Rank-8 LoRA on qkv/proj + head trained \\
 & Pretrained & External weights (trained on NIGHTS-20k) \\
 & Scratch & Retrained on the stated dataset from scratch \\
\midrule
Head & CLS-cos & Cosine distance on CLS token \\
 & Tversky-frozen & Tversky contrast, $\alpha=\beta=0.5$ \\
 & Cog-v1 & Learned cross-attention + context + Tversky \\
 & Cog-v2 & Shared alignment + bottleneck + Tversky \\
 & Cog & Frozen cosine align + GCM context + Tversky \\
 & Cog-strat & Cog with inverse-frequency sampling \\
 & TNN & Doumbouya 2025 differentiable TNN \\
 & Hinge-ens & Margin hinge loss over ensemble features \\
 & VLM-judge & Qwen3-VL logit ratio on A/B choice \\
\midrule
Dataset & NIGHTS-100k & Full 100k, 64{,}420 non-tie, 35{,}580 ties held out \\
 & NIGHTS-100k-unan & X-0 subset, $N{=}20{,}017$ \\
 & NIGHTS-100k-weak & 2-1/3-1/4-1 subset, $N{=}36{,}476$ \\
 & NIGHTS-20k (published) & Filtered unanimous release \\
\bottomrule
\end{tabular}
\end{table}

\clearpage
\noindent{\Large\bfseries B\quad Per-fold raw numbers}
\label{app:perfold}
\vspace{1em}

Table~\ref{tab:main} in the main body reports the mean and standard deviation across the five cross-validation folds. Table~\ref{tab:perfold} lists the raw per-fold accuracy for each of the twelve evaluated methods so that reviewers can verify variance estimates and check for outlier folds. Folds are defined by fixed seed 42 on the full 64{,}420-triplet non-tie corpus; the same fold indices are reused across every method, so per-fold differences reflect model behavior rather than data shuffling. No fold deviates from its method's mean by more than 1 percentage point for any of the trained methods, and the across-method ordering is stable on every fold taken in isolation.

\begin{table}[!ht]
\caption{Per-fold accuracies (\%) on the NIGHTS-100k non-tie corpus (5-fold cross-validation, seed 42). Each row is one method identified by the naming scheme of Table~\ref{tab:naming}. Columns are the five test folds; the same fold definitions are reused across all methods. Mean and standard deviation over the five folds are reported in Table~\ref{tab:main} of the main body.}
\label{tab:perfold}
\centering
\small
\begin{tabular}{lccccc}
\toprule
Method & Fold 1 & Fold 2 & Fold 3 & Fold 4 & Fold 5 \\
\midrule
LPIPS-AlexNet-Pretrained & 65.4 & 65.5 & 65.3 & 65.3 & 65.0 \\
DINOv1-ZS-CLS-cos & 74.8 & 74.0 & 73.8 & 73.6 & 74.2 \\
DINOv2-ZS-CLS-cos & 70.8 & 71.0 & 70.8 & 70.3 & 70.0 \\
DINOv3-ZS-CLS-cos & 72.4 & 72.3 & 72.1 & 71.4 & 71.7 \\
DINOv1-LoRA-Cog & 74.8 & 75.6 & 75.0 & 75.3 & 74.9 \\
DINOv2-LoRA-Cog & 77.4 & 77.5 & 76.2 & 77.4 & 77.1 \\
DINOv2-LoRA-Cos & 78.7 & 78.4 & 77.6 & 78.3 & 78.8 \\
DINOv2-LoRA-Cog-strat & 76.4 & 76.4 & 76.2 & 75.8 & 77.2 \\
DINOv2-LoRA-TNN & 77.2 & 77.5 & 76.7 & 77.4 & 78.1 \\
DINOv3-LoRA-Cog & 78.7 & 78.9 & 78.3 & 78.4 & 79.1 \\
DreamSim-Ens-Pretrained & 80.5 & 80.9 & 80.4 & 80.9 & 80.8 \\
DreamSim-Ens-Scratch & 79.4 & 79.7 & 77.9 & 78.4 & 78.2 \\
\bottomrule
\end{tabular}
\end{table}

\clearpage
\noindent{\Large\bfseries C\quad Full Cohen's $\kappa$ matrix}
\label{app:kappa}
\vspace{1em}

Cohen's $\kappa$ quantifies the agreement between two classifiers after correcting for the agreement expected by chance. We compute $\kappa$ pairwise between eleven of the twelve methods in Table~\ref{tab:main} (the TNN head is omitted) on the pooled 64{,}420 non-tie test predictions (every triplet appears in exactly one fold's test split, so the pooled set is a disjoint union of the five held-out folds). Values near 1 indicate redundant methods that would gain little from ensembling; values near 0.4--0.6 indicate methods that make different errors and are candidates for ensembling. In Figure~\ref{fig:kappa} the three DINO zero-shot entries cluster tightly ($\kappa>0.63$), as do the five LoRA-adapted DINO variants ($\kappa>0.6$), confirming that the DINO family makes correlated errors. LPIPS sits further out ($\kappa=0.25\text{--}0.36$ with any DINO-family method), which is consistent with its different underlying feature basis. DreamSim-Pretrained and DreamSim-Scratch reach $\kappa=0.75$, the highest pairwise agreement in the matrix: retraining the DreamSim ensemble on NIGHTS-100k does not shift its error structure, it only shifts the overall accuracy by 2 pp. The scratch model is structurally the same predictor, trained on a cleaner subset, that fails on the same ambiguous triplets.

\begin{figure}[!ht]
\centering
\includegraphics[width=0.55\linewidth]{figures/kappa_matrix.pdf}
\caption{Cohen's $\kappa$ heatmap across eleven of the twelve methods (TNN omitted) on NIGHTS-100k non-tie predictions ($N=64{,}420$). Each cell is the $\kappa$ between the row and column method's binary left/right predictions on the same set of triplets. Diagonal is 1 by construction. $\kappa=1$ is perfect agreement; $\kappa=0$ is chance. Two blocks stand out: the DINO-LoRA block (within-block $\kappa>0.6$), and the two DreamSim variants that agree with each other at $\kappa=0.75$ but with LPIPS at only 0.25--0.36.}
\label{fig:kappa}
\end{figure}

\paragraph{Is $\kappa$ sufficient when labels have limited reproducibility?} One reviewer concern is that pooled $\kappa$ mixes agreement on reproducible-label triplets (X-0, ceiling 100\%) with agreement on low-reproducibility triplets (2-1, ceiling 33\%). If two methods ``agree'' on a 2-1 triplet, are they agreeing on shared signal or shared noise? We address this by stratifying $\kappa$ by agreement pattern: computing $\kappa$ separately on the X-0 (unanimous), X-1 (strong, $X\geq 5$), and weak (2-1/3-1/4-1) subsets. Table~\ref{tab:kappa-strat} shows that the ranking of redundancy is preserved across all strata, with magnitudes declining in the expected way on weak-majority triplets where label noise is higher.

\begin{table}[!ht]
\caption{Cohen's $\kappa$ stratified by human-agreement pattern on the pooled NIGHTS-100k non-tie predictions. The ranking of redundancy (DreamSim-Pretrained $\leftrightarrow$ DreamSim-Scratch highest, LPIPS $\leftrightarrow$ DreamSim lowest) is identical at every stratum. On the unanimous subset, where labels have 100\% reproducibility ceiling, the DreamSim-Pretrained $\leftrightarrow$ DreamSim-Scratch agreement rises to $\kappa=0.90$, the strongest piece of evidence that retraining the ensemble on NIGHTS-100k does not alter its decision surface.}
\label{tab:kappa-strat}
\centering
\small
\begin{tabular}{lcccc}
\toprule
Pair & Unanimous X-0 & Strong X-1 & Weak (2-1/3-1/4-1) & Pooled \\
     & $(N=20{,}005)$ & $(N=7{,}922)$ & $(N=36{,}476)$ & $(N=64{,}420)$ \\
\midrule
DreamSim-Pretrained $\leftrightarrow$ DreamSim-Scratch & 0.90 & 0.76 & 0.68 & 0.75 \\
DreamSim-Pretrained $\leftrightarrow$ DINOv3-LoRA-Cog & 0.89 & 0.69 & 0.58 & 0.69 \\
DINOv2-LoRA-Cos $\leftrightarrow$ DINOv3-LoRA-Cog    & 0.85 & 0.64 & 0.55 & 0.67 \\
DINOv3-ZS $\leftrightarrow$ DreamSim-Pretrained         & 0.72 & 0.52 & 0.44 & 0.54 \\
LPIPS $\leftrightarrow$ DreamSim-Pretrained              & 0.52 & 0.34 & 0.27 & 0.30 \\
\bottomrule
\end{tabular}
\end{table}

The pooled $\kappa$ in Figure~\ref{fig:kappa} and in the main body is therefore a weighted-average over regimes with different label reliability. The stratified view shows two things: first, the redundancy claim (DreamSim variants and DINO-LoRA pairs are structurally similar classifiers, LPIPS is independent) holds on the subset where labels are authoritative; second, the lower $\kappa$ magnitudes on weak-majority patterns partly reflect the reduced maximum $\kappa$ achievable when classifier outputs are compared against noisy labels. Both observations are consistent with the Section~3 reproducibility-ceiling analysis: on weak-majority triplets there is less reproducible signal for any classifier (or pair of classifiers) to latch onto.

\clearpage
\noindent{\Large\bfseries D\quad Additional experiments}
\label{app:extra}
\vspace{1em}

This section documents experiments that did not make Table~\ref{tab:main} but were run in full: alternative similarity heads on frozen DINOv3 CLS embeddings, a multi-teacher distillation backbone (C-RADIOv4), a vision-language-model-as-judge (Qwen3-VL), ClickMe attention harmonization, and qualitative traces that make the agreement-stratified story visible on individual triplets. The BAPPS ceiling and paired significance are reported in the main body (Sections~3.4--3.7); their tables and interpretations appear there.

\subsection*{D.0\quad Harmonized DINOv3 (attention alignment on ClickMe)}

\citet{fel2022harmonization} propose training a vision backbone to match human attention maps collected via ClickMe~\citep{linsley2019clickme}. The hypothesis is that attention-aligned features yield similarity-aligned embeddings. We tested this on DINOv3 with the authors' pretrained ClickMe attention predictor: 5-fold CV on NIGHTS-100k non-tie (matched splits to Table~\ref{tab:main}) gives Harmonized-DINOv3 $= 64.0 \pm 0.2\%$, an 8 pp drop from DINOv3 zero-shot on the same folds (Table~\ref{tab:harm}). The near-zero across-fold variance indicates the drop is not noise.

The interpretation in the Discussion is the level-of-abstraction mismatch: ClickMe captures where humans look during categorization, not which compositional features drive perceived ensemble-image similarity. The result is consistent with attention-alignment being a useful signal for LPIPS-style pixel-distortion metrics but not for the mid/high-level discrimination NIGHTS requires.

\begin{table}[h]
\caption{Harmonized DINOv3 on NIGHTS-100k, 5-fold CV. Zero-shot DINOv3 reproduced from Table~\ref{tab:main} for direct comparison on the same test splits.}
\label{tab:harm}
\centering
\small
\begin{tabular}{lccccccc}
\toprule
Method & F1 & F2 & F3 & F4 & F5 & Mean $\pm$ Std & $\Delta$ \\
\midrule
DINOv3-ZS-CLS-cos            & 72.4 & 72.3 & 72.1 & 71.4 & 71.7 & $72.0 \pm 0.4$ & --     \\
Harmonized DINOv3            & 63.8 & 64.1 & 63.9 & 64.4 & 63.9 & $64.0 \pm 0.2$ & $-8.0$ \\
\bottomrule
\end{tabular}
\end{table}

\subsection*{D.1\quad Alternative similarity heads on frozen CLS}

The main body establishes that cosine distance on DINO CLS tokens is the strongest zero-shot signal on NIGHTS-100k. A natural follow-up question is whether any other similarity function on the same frozen embedding improves that number. We swept ten alternatives spanning basic distances, kernel methods, non-Euclidean geometries, distribution-based interpretations, and learned metric-learning heads, all evaluated on the same 5-fold cross-validation on NIGHTS-100k non-tie (Table~\ref{tab:head-sweep}).

\begin{table}[!ht]
\caption{Alternative similarity heads on frozen DINOv3 CLS embeddings, 5-fold CV on NIGHTS-100k non-tie ($N=64{,}420$). Cosine is the baseline from Table~\ref{tab:main}. ``Learned'' heads are trained with pairwise BCE for 10 epochs; all other heads are parameter-free. No alternative beats plain cosine; learned heads with more parameters do systematically worse. Retrieval-style losses did not help either (cosine-softmax-ranking and margin-triplet variants of cosine all remained within 0.1 pp of plain cosine on a fold-0 pilot).}
\label{tab:head-sweep}
\centering
\small
\begin{tabular}{llrr}
\toprule
Head & Type & Trainable params & Accuracy \\
\midrule
\textbf{Cosine} (baseline)                & parameter-free, angular & 0        & \textbf{72.01 $\pm$ 0.4} \\
Neg-dot product                           & parameter-free, inner   & 0        & 72.01 $\pm$ 0.3 \\
Arc-cosine kernel (order-1)               & Cho \& Saul 2009        & 0        & 72.11 $\pm$ 0.3 \\
Poincar\'e hyperbolic distance            & non-Euclidean           & 0        & 72.09 $\pm$ 0.3 \\
L2 Euclidean                              & parameter-free, metric  & 0        & 71.56 $\pm$ 0.3 \\
L1 Manhattan                              & parameter-free, metric  & 0        & 71.50 $\pm$ 0.2 \\
JSD on softmax(CLS)                       & distribution divergence & 0        & 71.37 $\pm$ 0.2 \\
\midrule
Mahalanobis diag (learned)                & learned metric          & 768      & 71.79 $\pm$ 0.2 \\
Mahalanobis rank-32 (learned)             & learned metric          & 24{,}576 & 70.26 $\pm$ 0.4 \\
Learned bilinear rank-32                  & learned metric          & 49{,}152 & 68.35 $\pm$ 0.3 \\
Siamese MLP on $[a, b, |a{-}b|, a{\odot}b]$ & learned similarity    & $\sim$800k & 68.39 $\pm$ 0.5 \\
\bottomrule
\end{tabular}
\end{table}

Three observations. First, the top four entries cluster within 0.1 pp of each other (cosine, neg-dot, arc-cosine, Poincar\'e). These differ only by a monotone transform of the inner product, which is why they cannot yield different binary preferences on the same triplets. Second, L2, L1, and JSD are all $\sim$0.5--0.7 pp below cosine: once the head stops being purely angular, the signal degrades. Third, learned heads with more parameters get strictly worse: adding 24{,}576 parameters to a Mahalanobis metric loses 1.7 pp, and an 800k-parameter Siamese MLP loses 3.6 pp despite having full access to the embedding pair. The learned heads are overfitting the fold-specific label-noise pattern documented in Section~3, not finding a better similarity geometry. On frozen DINO CLS embeddings, cosine is the best head we tested: every alternative we tried, parameter-free or learned, did worse.

\subsection*{D.2\quad Multi-teacher distillation (negative result)}

C-RADIOv4-SO400M~\citep{conwell2024large} distills DINOv3, SigLIP2, and SAM3 into a single ViT. We hypothesized that a model with stronger semantic, linguistic, and segmentation priors would produce better perceptual similarity. We applied rank-8 LoRA to qkv/proj and a cosine head on the summary token and trained for one epoch on fold 0 (approximately 28 hours at the model's required 1008$\times$1008 input resolution for its RoPE embeddings).

Epoch-1 test accuracy on fold 0 was 58.2\%, below every DINO backbone in Table~\ref{tab:main} by 13 pp or more, and below LPIPS by 7 pp. We stopped after one epoch. The model's pretraining objectives mix spatial, semantic, and linguistic signals that appear to dilute the fine-grained synthetic-image similarity signal NIGHTS requires. Training at 224$\times$224 is not possible without retraining the RoPE embeddings, and the full-resolution compute cost makes multi-fold cross-validation impractical for this backbone at our compute budget.

\subsection*{D.2\quad VLM-as-judge (negative result)}

Modern vision-language models can describe image content at the scene level. We tested whether a VLM can directly answer the 2AFC question. Using Qwen3-VL-4B-Instruct on a 100-triplet sample from fold 0, we constructed a prompt showing the reference and two candidates with randomized A/B order, and extracted the logit ratio of the A vs.\ B tokens at the first generated position. We applied position randomization to remove A/B bias.

\begin{table}[!ht]
\caption{Qwen3-VL-4B-Instruct on a 100-triplet sample of NIGHTS-100k fold 0.}
\label{tab:qwen}
\centering
\small
\begin{tabular}{lcc}
\toprule
Triplet subset & N & Accuracy \\
\midrule
Unanimous (X-0)            & 42 & 92.9\% \\
Strong-majority (X-1, $X{\geq}5$) & 17 & 76.5\% \\
Weak-majority (2-1/3-1/4-1) & 41 & 56.1\% \\
\midrule
All                        & 100 & 63.0\% \\
\bottomrule
\end{tabular}
\end{table}

The VLM can reliably call unanimous triplets but falls to near chance on weak-majority pairs. Scene-level semantic reasoning fails on the subtle perceptual discrimination NIGHTS is built to test: the model can tell that two images both depict a dog but not which dog is structurally closer to the reference.

\clearpage
\noindent{\Large\bfseries D.4\quad Model confidence versus the human vote split}
\label{app:confidence}
\vspace{1em}

Method and headline numbers are in \S\ref{sec:confidence}. Figure~\ref{fig:confcal} reports the full picture: confidence by agreement pattern against the confidence the votes themselves support, a reliability curve of calibrated $P(\text{left})$ against the observed vote fraction, and mean confidence on the 35{,}580 split-vote triplets. Figure~\ref{fig:confcum} repeats the measurement cumulatively, starting from the published 20k and relaxing the agreement filter one vote pattern at a time out to the full 100k. The confidence the votes support falls from 1.000 to 0.414 across that sweep; DreamSim-Ens-Pretrained falls from 0.790 to 0.563, and the two curves cross near 70{,}000 triplets.

\begin{figure}[h]
\centering
\includegraphics[width=\linewidth]{figures/confidence_calibration.pdf}
\caption{Model confidence against the human vote split on NIGHTS-100k. (A) confidence by agreement pattern; black is the confidence the observed vote split supports. (B) reliability against the observed vote fraction; legend gives ECE against the vote fraction and against the binarized majority label. (C) mean confidence on the 35{,}580 triplets where humans split 50/50. Gray bars are LoRA-trained heads reporting their native sigmoid.}
\label{fig:confcal}
\end{figure}

\begin{figure}[h]
\centering
\includegraphics[width=\linewidth]{figures/confidence_by_agreement.pdf}
\caption{The same measurement on the agreement split used throughout the paper. (A) the 17 raw vote patterns on the majority-vote-fraction axis of Figure~\ref{fig:agreement}; the unanimous block is the published 20k. (B) cumulative corpus, relaxing the filter one pattern at a time from the published 20k to the full 100k.}
\label{fig:confcum}
\end{figure}

\clearpage
\noindent{\Large\bfseries D.5\quad Per-axis stimulus manipulation results}
\label{app:axes_detail}
\vspace{1em}

These are the per-axis results of the Figure-6-style replica described in \S\ref{sec:shape_axes}. We report them here rather than in the main body because NIGHTS provides no human judgments on manipulated stimuli, and we collected none: the orientation, color, shading and blur axes therefore measure model invariance only and cannot be set against human behavior. The foreground/background axis is treated separately in the main text, where IN-9 and ObjectNet supply an external anchor.

\begin{figure}[h]
\centering
\includegraphics[width=0.88\linewidth]{figures/fig6_ablation_heatmap.pdf}
\caption{DreamSim-style Figure-6 ablation replica. Accuracy drop (pp) relative to clean for each backbone $\times$ ablation. Bluer is more invariant; redder is more fragile. DINOv3 zero-shot is the most invariant overall; DreamSim-Ens-Pretrained and our own DINOv2+LoRA are the most fragile despite the best clean accuracy. 2000 unanimous NIGHTS triplets.}
\label{fig:fig6_replica}
\end{figure}

\begin{enumerate}[leftmargin=*,itemsep=1pt,topsep=1pt]
\item \textbf{DINOv3 was the most invariant zero-shot backbone} on every axis except background removal. Rotation drop was $\leq 0.7$ pp, grayscale drop was only 0.3 pp, mild blur actually \emph{helped} (+1.25 pp at $\sigma=4$), and the foreground-only drop was 1.1 pp. Its clean accuracy of 85.4\% was below DINOv1's 88.5\% and DreamSim-Ens-Pretrained's 93.8\%, but its invariance profile was cleaner.
\item \textbf{DreamSim-Ens-Pretrained traded invariance for clean-image accuracy.} It had the best clean score in the table ($93.8\%$) but the worst or near-worst drop on rotation ($-4.3$), color ($-5.7$), heavy blur ($-9.6$), foreground ($-9.05$), and background ($-14.5$). Margin-hinge LoRA training narrows the feature geometry around the triplet distribution, which produces the best in-distribution score at the cost of cross-manipulation stability.
\end{enumerate}

Both observations are statements about models. Neither licenses a claim about the difference between model error and human error on these stimuli, and the revision no longer makes one.

\clearpage
\noindent{\Large\bfseries E\quad Training details}
\label{app:training}
\vspace{1em}

This section documents the full training recipe for reproducibility. All trained methods in Table~\ref{tab:main} use the same pipeline except where explicitly noted; DreamSim-Scratch uses the authors' published training code verbatim and is documented separately at the end.

\paragraph{Preprocessing.} Images are decoded from the NIGHTS release JPEGs, resized to $224\times 224$ with bicubic interpolation, and center-cropped. Normalization follows the backbone: ImageNet statistics (mean $[0.485, 0.456, 0.406]$, std $[0.229, 0.224, 0.225]$) for DINOv1, DINOv2, and DINOv3; $0.5/0.5$ for C-RADIOv4; the processor-supplied normalization for DreamSim and Qwen3-VL. No data augmentation is used at training time: the triplet task is orientation-sensitive, and the reference/candidate pairing is the supervision signal.

\paragraph{LoRA configuration.} LoRA adapters have rank 8, $\alpha = 16$, and dropout 0.05, and are inserted on the \texttt{qkv} and \texttt{proj} linear layers of every transformer block. Layer-norm parameters of the backbone are also unfrozen. For DINOv2 ViT-B/14 with 12 blocks this yields 442{,}368 LoRA parameters plus 1{,}536 LayerNorm parameters, which is approximately 0.5\% of the 86M backbone. The DreamSim ensemble retraining uses rank 16 on the same modules to match the authors' released configuration.

\paragraph{Optimizer.} Schedule-Free AdamW \citep{defazio2024schedulefree} with $\beta_1 = 0.9$, $\beta_2 = 0.999$, weight decay $10^{-2}$, 100 linear warm-up steps, and a constant learning rate thereafter (no cosine schedule). Separate learning rates are used: $10^{-4}$ for LoRA and layer-norm parameters, $2\times 10^{-4}$ for the head. Gradients are clipped to a global norm of 1.0 before each optimizer step.

\paragraph{Batch and compute.} Batch size is 32 triplets per step (each triplet contributes three images to the forward pass). Training runs for 10 epochs over the fold's training split. Compute is a shared cluster of NVIDIA RTX A6000 GPUs (48\,GB). A full 5-fold run of DINOv2-LoRA-Cog takes approximately 3--5 hours when five GPUs run concurrently, one per fold.

\paragraph{Checkpoint selection.} A held-out validation split of 10\% of each fold's training data selects the best-of-10-epochs checkpoint by validation accuracy; the selected checkpoint is evaluated once on the fold's test split, and that single number appears in Table~\ref{tab:perfold}. Per-triplet test-set predictions are persisted for every (method, fold) pair so that the Cohen's $\kappa$ matrix, McNemar tests, and error-distribution analyses can be reproduced without re-running training.

\paragraph{DreamSim retraining from scratch.} We run the authors' published training code unchanged with the released ensemble configuration: DINO+CLIP+OpenCLIP, rank-16 LoRA, margin-0.05 hinge loss, batch size 16, 6 epochs. The release filters training triplets to those with at least 6 votes, so the scratch model sees only strong-majority triplets during training; we evaluate on the same full non-tie test fold as every other method in Table~\ref{tab:main}. This mismatch is the point of the comparison, not an oversight: it isolates what the authors' training recipe would produce when the 100k distribution is available.

\clearpage
\noindent{\Large\bfseries F\quad Shape-only class separability and retrieval}
\label{app:shape_only}
\vspace{1em}

\begin{table}[h]
\caption{Shape-only class separability. Leave-one-out nearest-prototype accuracy on the Geirhos filled-silhouette and edge stimuli. Chance is 6.25\% (16 classes). Referenced from main body \S 5.10.}
\label{tab:shape_only}
\centering
\small
\begin{tabular}{lcccc}
\toprule
Backbone & Silh 1-NN & Silh proto & Edges 1-NN & Edges proto \\
\midrule
ResNet-50        & 77.5 & 81.3 & 63.1 & 67.5 \\
DINOv1           & 71.3 & 78.1 & 66.3 & 68.1 \\
CLIP-ViT-B/16    & 73.1 & 78.8 & 79.4 & 81.3 \\
DINOv2           & 81.3 & 86.3 & 71.3 & 80.0 \\
DINOv3           & 82.5 & 87.5 & 86.9 & 91.3 \\
ConvNeXt-Large   & 84.4 & 87.5 & 90.0 & 93.8 \\
\bottomrule
\end{tabular}
\end{table}

\begin{figure}[h]
\centering
\includegraphics[width=\linewidth]{figures/silhouette_retrieval.png}
\caption{Shape-only nearest-neighbor retrieval. For four query silhouettes (bicycle, cat, elephant, keyboard), each backbone's top-3 nearest silhouettes by cosine distance in CLS space. Green outline: match; red: mismatch. Every backbone, including ResNet-50 with its near-zero cue-conflict shape-bias, retrieves predominantly within-class; shape representation is present in every backbone we tested.}
\label{fig:silh_retrieval}
\end{figure}

\begin{figure}[h]
\centering
\includegraphics[width=0.8\linewidth]{figures/cost_of_accuracy.pdf}
\caption{The cost of clean-image accuracy. Each point is a backbone; vertical axis is the mean magnitude of accuracy drop across the 11 Fig-6 ablation conditions. Zero-shot backbones (circles) cluster in the low-fragility region; LoRA-trained models and DreamSim-Pretrained (diamonds) shift up and right. Dashed line is a least-squares trend.}
\label{fig:cost_acc}
\end{figure}

\paragraph{Tie-based bias audit.} On the 7,116-triplet tie subset held out per fold, LoRA-trained methods predict left with frequency 49.5--51.1\% (no positional bias). Zero-shot DINOs show 49.6--50.8\% and DreamSim 51\%, also near chance. The one method with tie bias outside that range is Qwen3-VL-4B (55\% left, confidence 0.30), consistent with known VLM positional preferences. LoRA adaptation on NIGHTS-100k does not induce spurious tie preferences.

\clearpage
\noindent{\Large\bfseries G\quad Full 2AFC-vs-JND axis table}
\label{app:axes}
\vspace{1em}

\begin{table}[h]
\caption{Per-backbone 2AFC accuracy (Fig-6 clean, 2000 unanimous triplets) and JND AUC (608 pairs, same/different). Extended form of Figure~\ref{fig:tasks_decouple}. The 2AFC and JND rankings invert across DINO generations and more broadly across the backbone set.}
\label{tab:axes}
\centering
\small
\begin{tabular}{lccl}
\toprule
Backbone & 2AFC acc (\%) & JND AUC & Family \\
\midrule
DINOv1              & 88.45 & 0.680 & DINO self-distill \\
DINOv2              & 84.65 & 0.778 & DINO self-distill \\
DINOv3              & 85.35 & 0.781 & DINO self-distill \\
CLIP-ViT-B/16       & 80.60 & 0.745 & image-text \\
OpenCLIP-ViT-B/16   & 83.95 & 0.740 & image-text \\
ViT-B/16 & 86.90 & 0.721 & label-supervised ViT \\
ConvNeXt-Large      & 85.40 & 0.749 & modern CNN \\
ResNet-50           & 84.05 & 0.734 & classical CNN \\
ResNet-152          & 83.45 & 0.740 & classical CNN \\
BiT-M R101x3        & 83.50 & 0.690 & JFT-scale CNN \\
CORnet-S            & 86.60 & 0.701 & biologically-inspired \\
BagNet-33           & 84.20 & 0.666 & biologically-inspired \\
VOneResNet50$^\ddagger$ & 61.05 & 0.677 & V1 Gabor front-end, untrained \\
DreamSim-Ens-ZS     & 89.00 & 0.685 & concatenated ensemble, ZS \\
DINOv2 + LoRA       & 89.30 & 0.778 & LoRA-trained on NIGHTS \\
DreamSim-Pretrained & 93.75 & 0.680 & LoRA-trained ensemble \\
\bottomrule
\end{tabular}
\vspace{0.35em}

\raggedright
{\footnotesize $^\ddagger$VOneResNet50 is a biologically-motivated V1 Gabor front-end followed by an \emph{untrained} ResNet-50 head; its near-chance 61\% is expected and it is included as a biologically-anchored control, not as a competitive method.}
\end{table}

\clearpage
\noindent{\Large\bfseries H\quad BAPPS-vs-NIGHTS quadrant view}
\label{app:crossover}
\vspace{1em}

\begin{figure}[h]
\centering
\includegraphics[width=0.45\linewidth]{figures/bapps_vs_nights_crossover.pdf}
\caption{Quadrant view of Figure~\ref{fig:crossover}. Each benchmark's split-half noise floor is drawn as a dashed line. The top-right quadrant (above both noise floors) is the only region with a generalizable cross-level alignment claim; only DreamSim-Ens-Pretrained sits inside it, and only because it was supervised-trained on the NIGHTS-20k subset.}
\label{fig:crossover_quadrant}
\end{figure}

\begin{figure}[h]
\centering
\includegraphics[width=0.72\linewidth]{figures/bapps_vs_nights_spectrum.pdf}
\caption{Perceptual similarity is a spectrum. Each method's accuracy at low-level BAPPS (left) connects to its accuracy at mid-to-high-level NIGHTS-100k (right). Crossing lines show rank inversion across the spectrum: LPIPS is a BAPPS specialist that falls below the NIGHTS noise floor; DINOv2/v3 invert the pattern. DreamSim is the only method above both noise floors, and only after supervised NIGHTS-20k training.}
\label{fig:crossover}
\end{figure}

\begin{table}[h]
\caption{Split-half reproducibility ceiling for BAPPS 2AFC validation (5 votes per triplet). Overall 69.0\% matches NIGHTS-100k (70.1\%) within 1 pp.}
\label{tab:bapps-ceiling}
\centering
\small
\begin{tabular}{lcccccc}
\toprule
Category & CNN & Traditional & Super-res & Color & Frame-int & Deblur \\
\midrule
N        & 4{,}720 & 4{,}720 & 10{,}856 & 4{,}720 & 1{,}888 & 9{,}440 \\
Ceiling  & 83.2\% & 79.1\% & 70.0\% & 64.1\% & 62.8\% & 59.6\% \\
\bottomrule
\end{tabular}
\quad\textbf{All BAPPS: N=36{,}344, ceiling = 69.0\%.}
\end{table}

\begin{table}[h]
\caption{Accuracy on ImageNet-prompted and non-ImageNet-prompted subsets of NIGHTS-100k non-tie. $\Delta$ is ImageNet minus non-ImageNet.}
\label{tab:imnet}
\centering
\small
\begin{tabular}{lccc}
\toprule
Method & ImageNet & Non-ImageNet & $\Delta$ \\
\midrule
LPIPS-AlexNet-Pretrained       & 64.6 & 66.1 & $-1.5$ \\
DINOv1-ZS-CLS-cos              & 74.5 & 73.6 & $+0.9$ \\
DINOv3-ZS-CLS-cos              & 73.0 & 70.7 & $+2.3$ \\
DINOv2-LoRA-Cos                & 78.7 & 77.9 & $+0.7$ \\
DINOv3-LoRA-Cog             & 79.4 & 77.9 & $+1.5$ \\
DreamSim-Ens-Pretrained        & 81.0 & 80.3 & $+0.8$ \\
DreamSim-Ens-Scratch           & 79.0 & 78.3 & $+0.7$ \\
\bottomrule
\end{tabular}
\end{table}

\begin{table}[h]
\caption{Zero-shot 2AFC accuracy on BAPPS validation set, averaged over six distortion categories (traditional, cnn, color, deblur, frame interpolation, super-resolution).}
\label{tab:bapps}
\centering
\small
\begin{tabular}{lc}
\toprule
Method & BAPPS average \\
\midrule
LPIPS-AlexNet               & 72.9\% \\
DreamSim-Ens-Pretrained     & 73.6\% \\
DINOv1-ZS                   & 71.5\% \\
DINOv2-ZS                   & 67.3\% \\
DINOv3-ZS                   & 68.4\% \\
\bottomrule
\end{tabular}
\end{table}

\begin{figure}[h]
\centering
\includegraphics[width=0.82\linewidth]{figures/error_distribution.pdf}
\caption{Left: share of each method's NIGHTS-100k errors, stratified by the human-agreement pattern of the underlying triplet (weak 2-1/3-1/4-1, strong X-1 with $X\geq 5$, unanimous X-0). Right: fraction of each method's wrong predictions where at least one human voter agreed with the model. The best methods are almost never alone against the humans (DreamSim 96.6\%, DINOv3-LoRA 92.0\%); their errors live where humans themselves were split.}
\label{fig:errdist}
\end{figure}

\begin{figure}[h]
\centering
\includegraphics[width=0.78\linewidth]{figures/jnd_vs_2afc_by_tier.pdf}
\caption{JND and 2AFC labels on the 608 shared stimuli, stratified by 2AFC vote count. Green cells = tasks agree; red cells = tasks disagree. The JND-2AFC disagreement rate is approximately 50\% at every 2AFC unanimity level, and is highest (64\%) on the strongest-consensus 10-0 bucket: more 2AFC voter agreement does not reduce JND disagreement.}
\label{fig:jnd_2afc_tier}
\end{figure}

\clearpage
\noindent{\Large\bfseries I\quad Extended §4.7 figures: shape bias, hGRU collapse, LoRA trajectory, JND task inversion}
\label{app:shape_axes_extra}
\vspace{1em}

These five figures were referenced from the main body §4.7 but deferred here for length.

\begin{figure}[h]
\centering
\includegraphics[width=0.68\linewidth]{figures/shape_bias_cue_conflict.pdf}
\caption{Cue-conflict shape bias on the Geirhos 2018 stimuli. DINOv1 preferences are texture-dominated (10.3\%); DINOv2 (23.2\%) and DINOv3 (37.3\%) are more than three times more shape-biased. ResNet-50 is the most texture-dominant (3.75\%); ConvNeXt-Large (19.75\%) and CLIP-ViT-B/16 (29.83\%) sit between. Zero-shot, 1200 cue-conflict images, 16 classes. Humans score $\sim$95\% on the same stimuli \citep{geirhos2018imagenettrained}.}
\label{fig:shape_bias_cc}
\end{figure}

\begin{figure}[h]
\centering
\includegraphics[width=0.78\linewidth]{figures/hgru_spatial_sweep.pdf}
\caption{Radial-mask sensitivity sweep. Accuracy on 2000 unanimous triplets as a function of radial mask radius, keeping only the center disc (fg, left) or only the surround ring (bg, right). DINO backbones transition smoothly; ResNet-50+hGRU collapses to chance on any foreground radius in $[0.1, 0.5]$ and any background radius below $0.7$, revealing that the lateral-recurrence stage learns a globally coupled representation that fails under any spatial mask.}
\label{fig:hgru_sweep}
\end{figure}

\begin{figure}[h]
\centering
\includegraphics[width=0.92\linewidth]{figures/attention_zs_vs_lora.pdf}
\caption{Three high-vote NIGHTS unanimous triplets where DINOv2 + LoRA flips the majority-choice decision. Green border = human-majority candidate; red border = LoRA's (wrong) choice. Overlay: last-layer CLS attention. Text panels report the exact $\cos(\mathrm{ref}, \cdot)$ the models used. Two consistent signatures: (i) LoRA's flips land on $\cos$ margins as small as 0.005 (spoonbill: $0.905$ vs $0.910$), exactly the voter-noise-scale differences the reproducibility analysis of Section~\ref{sec:ceiling} warned about; (ii) attention geometry is nearly identical between ZS and LoRA, so the shortcut is encoded in the feature representation (sim(full, fg-only) drops consistently for LoRA) rather than in where the model looks. Pairwise similarities computed with DINOv2 + rank-8 LoRA + cosine at 224$\times$224.}
\label{fig:attn_zs_vs_lora}
\end{figure}

\begin{figure}[h]
\centering
\includegraphics[width=0.78\linewidth]{figures/lora_strength_sweep.pdf}
\caption{Controlled LoRA-strength sweep on a single backbone. DINOv2 + rank-8 LoRA trained for 0--8 epochs on NIGHTS-20k unanimous, each checkpoint evaluated on the 2000-triplet Fig-6 ablation suite. Gray dots: 16-backbone context from Figure~\ref{fig:cost_acc}. Phase 1 (epoch 0\,$\to$\,1) mirrors the cross-backbone invariance tax slope. Phase 2 (epochs 1--8) keeps climbing the fragility axis while clean accuracy reverts toward the zero-shot baseline. Within-backbone evidence that the invariance tax is causal, not a family-level confound.}
\label{fig:lora_strength}
\end{figure}

\begin{figure}[h]
\centering
\includegraphics[width=0.78\linewidth]{figures/tasks_decouple_scatter.pdf}
\caption{Perceptual alignment is not one axis. Each point is one backbone; horizontal axis is NIGHTS 2AFC accuracy on the 2000 unanimous-triplet subset; vertical axis is JND AUC on the 608-pair same/different corpus. The green arrows trace the DINOv1 $\to$ v2 $\to$ v3 progression: v1 wins 2AFC but loses JND, v2 and v3 win JND but lose 2AFC to v1. Same backbones, same image distribution (NIGHTS), different tasks produce different rankings.}
\label{fig:tasks_decouple}
\end{figure}


\clearpage
\noindent{\Large\bfseries J\quad Qualitative triplet traces}
\label{app:traces}
\vspace{1em}

The analyses in the main body are aggregate. This section shows the individual stimuli behind them, drawn unmodified from the released corpus. Figure~\ref{fig:traces_agreement} illustrates what the agreement patterns of \S\ref{sec:ceiling} look like on single triplets, one per pattern from 2-1 to 7-0. A 2-1 triplet carries a label a fresh panel of three annotators would reproduce a third of the time; a 7-0 triplet carries one that reproduces every time. Both contribute equally to a headline accuracy.

\begin{figure}[p]
\centering
\includegraphics[height=0.86\textheight]{figures/triplet_traces_by_agreement.png}
\caption{One NIGHTS-2AFC triplet per vote pattern. Reference, the candidate the human majority chose, and the alternative, with the vote split annotators produced. Prompts and triplet ids are given beside each row.}
\label{fig:traces_agreement}
\end{figure}

Figure~\ref{fig:traces_hard} shows the complementary case: triplets with 8 to 10 votes on one side, where the label is as stable as NIGHTS labels get, and where DreamSim-Ens-Pretrained still selects the other candidate. These are the errors the reproducibility ceiling does not explain, and they are 3.4\% of that method's errors (\S\ref{sec:errors}).

\begin{figure}[!ht]
\centering
\includegraphics[width=\linewidth]{figures/unanimous_hard_examples.png}
\caption{Unanimous triplets that the strongest method in Table~\ref{tab:main} gets wrong. Vote pattern and prompt are given under each panel.}
\label{fig:traces_hard}
\end{figure}

Figure~\ref{fig:traces_jnd} shows stimuli carrying both NIGHTS labels. All 608 shared stimuli come from the X-0 unanimous tier, so the 2AFC/JND comparison of \S\ref{sec:discussion} is a comparison on the easy subset rather than the hard one, and the two labels still contradict each other on 51.6\% of them.

\begin{figure}[!ht]
\centering
\includegraphics[width=0.72\linewidth]{figures/jnd_2afc_disagreement_examples.png}
\caption{Stimuli whose NIGHTS-JND label disagrees with a unanimous NIGHTS-2AFC label. Each panel is a reference and one candidate, with the zero-shot and LoRA cosine distances printed beneath.}
\label{fig:traces_jnd}
\end{figure}

\clearpage


% !TEX root = cosine_is_human.tex
\section*{NeurIPS Paper Checklist}

\begin{enumerate}

\item {\bf Claims}
    \item[] Question: Do the main claims made in the abstract and introduction accurately reflect the paper's contributions and scope?
    \item[] Answer: \answerYes{} % Replace by \answerYes{}, \answerNo{}, or \answerNA{}.
    \item[] Justification: The abstract and introduction state the thesis (cosine on a self-supervised backbone captures most of the reproducible perceptual signal NIGHTS measures) and the supporting evidence: a closed-form reproducibility ceiling, foreground-only fragility under fine-tuning, and 2AFC--JND label dissociation. All numerical claims appear in the body and are reproduced from the released per-fold predictions.
    \item[] Guidelines:
    \begin{itemize}
        \item The answer \answerNA{} means that the abstract and introduction do not include the claims made in the paper.
        \item The abstract and/or introduction should clearly state the claims made, including the contributions made in the paper and important assumptions and limitations. A \answerNo{} or \answerNA{} answer to this question will not be perceived well by the reviewers. 
        \item The claims made should match theoretical and experimental results, and reflect how much the results can be expected to generalize to other settings. 
        \item It is fine to include aspirational goals as motivation as long as it is clear that these goals are not attained by the paper. 
    \end{itemize}

\item {\bf Limitations}
    \item[] Question: Does the paper discuss the limitations of the work performed by the authors?
    \item[] Answer: \answerYes{} % Replace by \answerYes{}, \answerNo{}, or \answerNA{}.
    \item[] Justification: Section~\ref{sec:limitations} lists seven explicit caveats: NIGHTS-specific method comparisons, the scope of backbone coverage, fixed LoRA rank, the ambiguous-pair vs minority-bias interpretation of the agreement statistic, the qualitative scale of the JND audit, the within-population scope of the reproducibility ceiling, and the absence of per-rater data and demographics.
    \item[] Guidelines:
    \begin{itemize}
        \item The answer \answerNA{} means that the paper has no limitation while the answer \answerNo{} means that the paper has limitations, but those are not discussed in the paper. 
        \item The authors are encouraged to create a separate ``Limitations'' section in their paper.
        \item The paper should point out any strong assumptions and how robust the results are to violations of these assumptions (e.g., independence assumptions, noiseless settings, model well-specification, asymptotic approximations only holding locally). The authors should reflect on how these assumptions might be violated in practice and what the implications would be.
        \item The authors should reflect on the scope of the claims made, e.g., if the approach was only tested on a few datasets or with a few runs. In general, empirical results often depend on implicit assumptions, which should be articulated.
        \item The authors should reflect on the factors that influence the performance of the approach. For example, a facial recognition algorithm may perform poorly when image resolution is low or images are taken in low lighting. Or a speech-to-text system might not be used reliably to provide closed captions for online lectures because it fails to handle technical jargon.
        \item The authors should discuss the computational efficiency of the proposed algorithms and how they scale with dataset size.
        \item If applicable, the authors should discuss possible limitations of their approach to address problems of privacy and fairness.
        \item While the authors might fear that complete honesty about limitations might be used by reviewers as grounds for rejection, a worse outcome might be that reviewers discover limitations that aren't acknowledged in the paper. The authors should use their best judgment and recognize that individual actions in favor of transparency play an important role in developing norms that preserve the integrity of the community. Reviewers will be specifically instructed to not penalize honesty concerning limitations.
    \end{itemize}

\item {\bf Theory assumptions and proofs}
    \item[] Question: For each theoretical result, does the paper provide the full set of assumptions and a complete (and correct) proof?
    \item[] Answer: \answerNA{} % Replace by \answerYes{}, \answerNo{}, or \answerNA{}.
    \item[] Justification: The paper makes empirical claims about benchmark behaviour. The reproducibility ceiling in Section~\ref{sec:ceiling} is computed in closed form by resampling each triplet's votes; no formal theorems are stated.
    \item[] Guidelines:
    \begin{itemize}
        \item The answer \answerNA{} means that the paper does not include theoretical results. 
        \item All the theorems, formulas, and proofs in the paper should be numbered and cross-referenced.
        \item All assumptions should be clearly stated or referenced in the statement of any theorems.
        \item The proofs can either appear in the main paper or the supplemental material, but if they appear in the supplemental material, the authors are encouraged to provide a short proof sketch to provide intuition. 
        \item Inversely, any informal proof provided in the core of the paper should be complemented by formal proofs provided in appendix or supplemental material.
        \item Theorems and Lemmas that the proof relies upon should be properly referenced. 
    \end{itemize}

    \item {\bf Experimental result reproducibility}
    \item[] Question: Does the paper fully disclose all the information needed to reproduce the main experimental results of the paper to the extent that it affects the main claims and/or conclusions of the paper (regardless of whether the code and data are provided or not)?
    \item[] Answer: \answerYes{} % Replace by \answerYes{}, \answerNo{}, or \answerNA{}.
    \item[] Justification: All twelve methods use the same fold definitions (KFold, seed 42, sorted by triplet id); training hyperparameters appear in Section~3.3 and Appendix~E; per-triplet test predictions for every method and fold are in the released code package.
    \item[] Guidelines:
    \begin{itemize}
        \item The answer \answerNA{} means that the paper does not include experiments.
        \item If the paper includes experiments, a \answerNo{} answer to this question will not be perceived well by the reviewers: Making the paper reproducible is important, regardless of whether the code and data are provided or not.
        \item If the contribution is a dataset and\slash or model, the authors should describe the steps taken to make their results reproducible or verifiable. 
        \item Depending on the contribution, reproducibility can be accomplished in various ways. For example, if the contribution is a novel architecture, describing the architecture fully might suffice, or if the contribution is a specific model and empirical evaluation, it may be necessary to either make it possible for others to replicate the model with the same dataset, or provide access to the model. In general. releasing code and data is often one good way to accomplish this, but reproducibility can also be provided via detailed instructions for how to replicate the results, access to a hosted model (e.g., in the case of a large language model), releasing of a model checkpoint, or other means that are appropriate to the research performed.
        \item While NeurIPS does not require releasing code, the conference does require all submissions to provide some reasonable avenue for reproducibility, which may depend on the nature of the contribution. For example
        \begin{enumerate}
            \item If the contribution is primarily a new algorithm, the paper should make it clear how to reproduce that algorithm.
            \item If the contribution is primarily a new model architecture, the paper should describe the architecture clearly and fully.
            \item If the contribution is a new model (e.g., a large language model), then there should either be a way to access this model for reproducing the results or a way to reproduce the model (e.g., with an open-source dataset or instructions for how to construct the dataset).
            \item We recognize that reproducibility may be tricky in some cases, in which case authors are welcome to describe the particular way they provide for reproducibility. In the case of closed-source models, it may be that access to the model is limited in some way (e.g., to registered users), but it should be possible for other researchers to have some path to reproducing or verifying the results.
        \end{enumerate}
    \end{itemize}


\item {\bf Open access to data and code}
    \item[] Question: Does the paper provide open access to the data and code, with sufficient instructions to faithfully reproduce the main experimental results, as described in supplemental material?
    \item[] Answer: \answerYes{} % Replace by \answerYes{}, \answerNo{}, or \answerNA{}.
    \item[] Justification: The released code package contains the analysis pipeline, fold definitions, per-triplet predictions for all twelve methods, retrained checkpoints, and the standalone reproducibility-ceiling script. NIGHTS and BAPPS are pre-existing public datasets and are not redistributed; their original sources are cited.
    \item[] Guidelines:
    \begin{itemize}
        \item The answer \answerNA{} means that paper does not include experiments requiring code.
        \item Please see the NeurIPS code and data submission guidelines (\url{https://neurips.cc/public/guides/CodeSubmissionPolicy}) for more details.
        \item While we encourage the release of code and data, we understand that this might not be possible, so \answerNo{} is an acceptable answer. Papers cannot be rejected simply for not including code, unless this is central to the contribution (e.g., for a new open-source benchmark).
        \item The instructions should contain the exact command and environment needed to run to reproduce the results. See the NeurIPS code and data submission guidelines (\url{https://neurips.cc/public/guides/CodeSubmissionPolicy}) for more details.
        \item The authors should provide instructions on data access and preparation, including how to access the raw data, preprocessed data, intermediate data, and generated data, etc.
        \item The authors should provide scripts to reproduce all experimental results for the new proposed method and baselines. If only a subset of experiments are reproducible, they should state which ones are omitted from the script and why.
        \item At submission time, to preserve anonymity, the authors should release anonymized versions (if applicable).
        \item Providing as much information as possible in supplemental material (appended to the paper) is recommended, but including URLs to data and code is permitted.
    \end{itemize}


\item {\bf Experimental setting/details}
    \item[] Question: Does the paper specify all the training and test details (e.g., data splits, hyperparameters, how they were chosen, type of optimizer) necessary to understand the results?
    \item[] Answer: \answerYes{} % Replace by \answerYes{}, \answerNo{}, or \answerNA{}.
    \item[] Justification: Methods, hyperparameters, and dataset processing are documented in Section~3 (Method) and Appendix~E (Training details). DreamSim-Scratch uses the authors' published training code unchanged, also documented in Appendix~E.
    \item[] Guidelines:
    \begin{itemize}
        \item The answer \answerNA{} means that the paper does not include experiments.
        \item The experimental setting should be presented in the core of the paper to a level of detail that is necessary to appreciate the results and make sense of them.
        \item The full details can be provided either with the code, in appendix, or as supplemental material.
    \end{itemize}

\item {\bf Experiment statistical significance}
    \item[] Question: Does the paper report error bars suitably and correctly defined or other appropriate information about the statistical significance of the experiments?
    \item[] Answer: \answerYes{} % Replace by \answerYes{}, \answerNo{}, or \answerNA{}.
    \item[] Justification: Table~\ref{tab:main} reports 5-fold mean and standard deviation. We do not report pairwise significance tests: gaps inside the label-noise band quantified in \S\ref{sec:ceiling} are not separable from label noise regardless of significance.
    \item[] Guidelines:
    \begin{itemize}
        \item The answer \answerNA{} means that the paper does not include experiments.
        \item The authors should answer \answerYes{} if the results are accompanied by error bars, confidence intervals, or statistical significance tests, at least for the experiments that support the main claims of the paper.
        \item The factors of variability that the error bars are capturing should be clearly stated (for example, train/test split, initialization, random drawing of some parameter, or overall run with given experimental conditions).
        \item The method for calculating the error bars should be explained (closed form formula, call to a library function, bootstrap, etc.)
        \item The assumptions made should be given (e.g., Normally distributed errors).
        \item It should be clear whether the error bar is the standard deviation or the standard error of the mean.
        \item It is OK to report 1-sigma error bars, but one should state it. The authors should preferably report a 2-sigma error bar than state that they have a 96\% CI, if the hypothesis of Normality of errors is not verified.
        \item For asymmetric distributions, the authors should be careful not to show in tables or figures symmetric error bars that would yield results that are out of range (e.g., negative error rates).
        \item If error bars are reported in tables or plots, the authors should explain in the text how they were calculated and reference the corresponding figures or tables in the text.
    \end{itemize}

\item {\bf Experiments compute resources}
    \item[] Question: For each experiment, does the paper provide sufficient information on the computer resources (type of compute workers, memory, time of execution) needed to reproduce the experiments?
    \item[] Answer: \answerYes{} % Replace by \answerYes{}, \answerNo{}, or \answerNA{}.
    \item[] Justification: Each LoRA training run completes in roughly 8 GPU-hours on a single A6000. The full 5-fold cross-validation across all twelve methods uses approximately 400 A6000-hours; zero-shot evaluation is approximately 30 GPU-minutes per backbone.
    \item[] Guidelines:
    \begin{itemize}
        \item The answer \answerNA{} means that the paper does not include experiments.
        \item The paper should indicate the type of compute workers CPU or GPU, internal cluster, or cloud provider, including relevant memory and storage.
        \item The paper should provide the amount of compute required for each of the individual experimental runs as well as estimate the total compute. 
        \item The paper should disclose whether the full research project required more compute than the experiments reported in the paper (e.g., preliminary or failed experiments that didn't make it into the paper). 
    \end{itemize}
    
\item {\bf Code of ethics}
    \item[] Question: Does the research conducted in the paper conform, in every respect, with the NeurIPS Code of Ethics \url{https://neurips.cc/public/EthicsGuidelines}?
    \item[] Answer: \answerYes{} % Replace by \answerYes{}, \answerNo{}, or \answerNA{}.
    \item[] Justification: The paper conforms to the NeurIPS Code of Ethics. Limitations and broader impact are explicitly discussed in Section~\ref{sec:limitations} and the Broader Impact section.
    \item[] Guidelines:
    \begin{itemize}
        \item The answer \answerNA{} means that the authors have not reviewed the NeurIPS Code of Ethics.
        \item If the authors answer \answerNo, they should explain the special circumstances that require a deviation from the Code of Ethics.
        \item The authors should make sure to preserve anonymity (e.g., if there is a special consideration due to laws or regulations in their jurisdiction).
    \end{itemize}


\item {\bf Broader impacts}
    \item[] Question: Does the paper discuss both potential positive societal impacts and negative societal impacts of the work performed?
    \item[] Answer: \answerYes{} % Replace by \answerYes{}, \answerNo{}, or \answerNA{}.
    \item[] Justification: Discussed in the Broader Impact and Reproducibility section. The principal effect is improved research hygiene around a widely used benchmark; we identify no direct negative societal impact.
    \item[] Guidelines:
    \begin{itemize}
        \item The answer \answerNA{} means that there is no societal impact of the work performed.
        \item If the authors answer \answerNA{} or \answerNo, they should explain why their work has no societal impact or why the paper does not address societal impact.
        \item Examples of negative societal impacts include potential malicious or unintended uses (e.g., disinformation, generating fake profiles, surveillance), fairness considerations (e.g., deployment of technologies that could make decisions that unfairly impact specific groups), privacy considerations, and security considerations.
        \item The conference expects that many papers will be foundational research and not tied to particular applications, let alone deployments. However, if there is a direct path to any negative applications, the authors should point it out. For example, it is legitimate to point out that an improvement in the quality of generative models could be used to generate Deepfakes for disinformation. On the other hand, it is not needed to point out that a generic algorithm for optimizing neural networks could enable people to train models that generate Deepfakes faster.
        \item The authors should consider possible harms that could arise when the technology is being used as intended and functioning correctly, harms that could arise when the technology is being used as intended but gives incorrect results, and harms following from (intentional or unintentional) misuse of the technology.
        \item If there are negative societal impacts, the authors could also discuss possible mitigation strategies (e.g., gated release of models, providing defenses in addition to attacks, mechanisms for monitoring misuse, mechanisms to monitor how a system learns from feedback over time, improving the efficiency and accessibility of ML).
    \end{itemize}
    
\item {\bf Safeguards}
    \item[] Question: Does the paper describe safeguards that have been put in place for responsible release of data or models that have a high risk for misuse (e.g., pre-trained language models, image generators, or scraped datasets)?
    \item[] Answer: \answerNA{} % Replace by \answerYes{}, \answerNo{}, or \answerNA{}.
    \item[] Justification: We release analysis code, vote-count statistics, and similarity-metric checkpoints with no high-risk capabilities (no generative model weights, no demographic-classification models, no surveillance-relevant artefacts).
    \item[] Guidelines:
    \begin{itemize}
        \item The answer \answerNA{} means that the paper poses no such risks.
        \item Released models that have a high risk for misuse or dual-use should be released with necessary safeguards to allow for controlled use of the model, for example by requiring that users adhere to usage guidelines or restrictions to access the model or implementing safety filters. 
        \item Datasets that have been scraped from the Internet could pose safety risks. The authors should describe how they avoided releasing unsafe images.
        \item We recognize that providing effective safeguards is challenging, and many papers do not require this, but we encourage authors to take this into account and make a best faith effort.
    \end{itemize}

\item {\bf Licenses for existing assets}
    \item[] Question: Are the creators or original owners of assets (e.g., code, data, models), used in the paper, properly credited and are the license and terms of use explicitly mentioned and properly respected?
    \item[] Answer: \answerYes{} % Replace by \answerYes{}, \answerNo{}, or \answerNA{}.
    \item[] Justification: NIGHTS~\citep{fu2023dreamsim} and BAPPS~\citep{zhang2018lpips} are used under their published research-use terms. DINOv1/v2/v3 backbones are used under their respective public releases. All assets are credited in the bibliography.
    \item[] Guidelines:
    \begin{itemize}
        \item The answer \answerNA{} means that the paper does not use existing assets.
        \item The authors should cite the original paper that produced the code package or dataset.
        \item The authors should state which version of the asset is used and, if possible, include a URL.
        \item The name of the license (e.g., CC-BY 4.0) should be included for each asset.
        \item For scraped data from a particular source (e.g., website), the copyright and terms of service of that source should be provided.
        \item If assets are released, the license, copyright information, and terms of use in the package should be provided. For popular datasets, \url{paperswithcode.com/datasets} has curated licenses for some datasets. Their licensing guide can help determine the license of a dataset.
        \item For existing datasets that are re-packaged, both the original license and the license of the derived asset (if it has changed) should be provided.
        \item If this information is not available online, the authors are encouraged to reach out to the asset's creators.
    \end{itemize}

\item {\bf New assets}
    \item[] Question: Are new assets introduced in the paper well documented and is the documentation provided alongside the assets?
    \item[] Answer: \answerYes{} % Replace by \answerYes{}, \answerNo{}, or \answerNA{}.
    \item[] Justification: The released code package includes a README, requirements.txt, MIT LICENSE, and per-method documentation. Fold definitions and per-triplet predictions are documented inline in the analysis package.
    \item[] Guidelines:
    \begin{itemize}
        \item The answer \answerNA{} means that the paper does not release new assets.
        \item Researchers should communicate the details of the dataset\slash code\slash model as part of their submissions via structured templates. This includes details about training, license, limitations, etc. 
        \item The paper should discuss whether and how consent was obtained from people whose asset is used.
        \item At submission time, remember to anonymize your assets (if applicable). You can either create an anonymized URL or include an anonymized zip file.
    \end{itemize}

\item {\bf Crowdsourcing and research with human subjects}
    \item[] Question: For crowdsourcing experiments and research with human subjects, does the paper include the full text of instructions given to participants and screenshots, if applicable, as well as details about compensation (if any)? 
    \item[] Answer: \answerNA{} % Replace by \answerYes{}, \answerNo{}, or \answerNA{}.
    \item[] Justification: We do not run new human-subjects studies. All similarity labels are re-used from the public NIGHTS and BAPPS releases.
    \item[] Guidelines:
    \begin{itemize}
        \item The answer \answerNA{} means that the paper does not involve crowdsourcing nor research with human subjects.
        \item Including this information in the supplemental material is fine, but if the main contribution of the paper involves human subjects, then as much detail as possible should be included in the main paper. 
        \item According to the NeurIPS Code of Ethics, workers involved in data collection, curation, or other labor should be paid at least the minimum wage in the country of the data collector. 
    \end{itemize}

\item {\bf Institutional review board (IRB) approvals or equivalent for research with human subjects}
    \item[] Question: Does the paper describe potential risks incurred by study participants, whether such risks were disclosed to the subjects, and whether Institutional Review Board (IRB) approvals (or an equivalent approval/review based on the requirements of your country or institution) were obtained?
    \item[] Answer: \answerNA{} % Replace by \answerYes{}, \answerNo{}, or \answerNA{}.
    \item[] Justification: No new human-subjects research was conducted. We re-analyse pre-existing public annotation data.
    \item[] Guidelines:
    \begin{itemize}
        \item The answer \answerNA{} means that the paper does not involve crowdsourcing nor research with human subjects.
        \item Depending on the country in which research is conducted, IRB approval (or equivalent) may be required for any human subjects research. If you obtained IRB approval, you should clearly state this in the paper. 
        \item We recognize that the procedures for this may vary significantly between institutions and locations, and we expect authors to adhere to the NeurIPS Code of Ethics and the guidelines for their institution. 
        \item For initial submissions, do not include any information that would break anonymity (if applicable), such as the institution conducting the review.
    \end{itemize}

\item {\bf Declaration of LLM usage}
    \item[] Question: Does the paper describe the usage of LLMs if it is an important, original, or non-standard component of the core methods in this research? Note that if the LLM is used only for writing, editing, or formatting purposes and does \emph{not} impact the core methodology, scientific rigor, or originality of the research, declaration is not required.
    %this research? 
    \item[] Answer: \answerYes{} % Replace by \answerYes{}, \answerNo{}, or \answerNA{}.
    \item[] Justification: A large language model was used as a writing assistant on parts of the manuscript and as a coding assistant during analysis script development. The LLM did not contribute to research ideation, experiment design, or data interpretation; all empirical results come from the released code and predictions and are independently reproducible.
    \item[] Guidelines:
    \begin{itemize}
        \item The answer \answerNA{} means that the core method development in this research does not involve LLMs as any important, original, or non-standard components.
        \item Please refer to our LLM policy in the NeurIPS handbook for what should or should not be described.
    \end{itemize}

\end{enumerate}

\end{document}
